
We recall the background material for setting up the algebraic category $\SSBim$ of singular Soergel bimodules from a given realization. We then recall the categorification results and the Williamson-Soergel Hom formula, and then show that these results are valid for the deformed affine type A realization we are working with, under suitable assumptions on the ground field. We also briefly discuss a reformulation of the classical Satake isomorphism, which is categorified by the the functor $\GS_\zeta$ constructed in \cref{subsec the functor}.

%============================================================
\subsection{Deformed affine realizations}
\label{subsec- realizations}
%============================================================

% \BE{OOOOOH. Ok. Let's change all the $V$ to $\Lambda$ in this section, including the notation $V_A$. Then no conflict with geordie...}

\begin{defn} (See \cite[Definition 3.1]{EWGr4sb}) Let $(W,S)$ be a Coxeter system and $\Bbbk$ a commutative domain. A \emph{realization} of $W$ over $\Bbbk$ is the data of a free
$\Bbbk$ module $\Lambda$ with a $W$-action, together with a set of roots $\{\al_s\} \subset \Lambda$ and coroots $\{\al_s^\vee\} \subset \Lambda^*=\Hom_{\kk}(V,\kk)$ indexed by $s \in S$. We require that the action of
$s \in S$ on $\Lambda$ is given by the formula \[ s(v) = v - \langle \al_s^\vee, v \rangle \al_s. \] We also require $\langle \al_s^\vee, \al_s \rangle = 2$ for all $s$, and $\langle
\al_s^\vee, \al_t \rangle = 0$ whenever $m_{st} = 2$. The \emph{Cartan matrix} of the realization is the $S \times S$ matrix with entries $\langle \al_s^\vee, \al_t \rangle$.
\end{defn}

Given a Coxeter system $(W,S)$, each subset $I \subset S$ defines a parabolic subgroup $W_I$ of $W$, and $(W_I,I)$ itself is a Coxeter system. 
Given a realization $(\Lambda,\Lambda^*,\{\al_s\}_{s\in S},\{\al_s^\vee\}_{s \in S})$ of $(W,S)$, we can then obtain the realization $(\Lambda,\Lambda^*, \set{\al_s}_{s \in I}, \set{\al_s^\vee}_{s\in I})$ of $(W_I,I)$, remembering only the roots and coroots for simple reflections in $I$, and $W_I$ acts on $\Lambda$ by restriction.

For a more detailed exposition on realizations, see \cite[Section 5.7]{EMTW}.

\begin{example}
    For $(W,S)=(S_n, \set{s_1, s_2, \ldots, s_{n-1}})$, we have the permutation realization (over $\kk$) with $\Lambda=\bigoplus_{i=1}^n\kk y_i$. 
    The roots are $ \al_{s_i} := y_i-y_{i+1}$ and $\al_{s_i}^\vee:= y_i^*-y_{i+1}^*$, where $\{y_i^*\}$ is the dual basis to $\{y_i\}$. This induces the usual permutation action on the set $\{y_i\}$.
\end{example}

We shall primarily focus on the affine Weyl group $W_{\aff}$ in Type A, associated to the symmetric group $S_n$. Viewing $W_{\aff}$ as a Coxeter group, it is generated by a set of simple reflections $S = \{s_i\}$ indexed by $\Om = \ZZ/n\ZZ$. We have $m_{s_i s_j} = 3$ if $j = i \pm 1$, and $m_{s_i
s_j} = 2$ otherwise. In this paper all indices will be considered modulo $n$. We say that $\{s_1, \ldots, s_{n-1}\}$ generate the finite Weyl group $S_n$, while $s_0$ is the \emph{affine reflection}. Because it leaves out $s_0$, we call the parabolic subgroup $(S_n,\{s_1, \ldots, s_{n-1}\})$ by the name $W_{\hh{0}}$. 

\begin{example}\cite[Definition 2.8, with $z$ replaced by $\zeta$]{EJY1} \label{defn:Vz} Let $\zeta$ be a formal variable, let $\kk'$ be a commutative domain, and let $\kk = \AC_{\zeta} := \kk'[\zeta,\zeta^{-1}]$. Let $\Lambda$ be the free $\AC_\zeta$-module $\AC_\zeta^{\oplus n}$ with basis $\{x_i\}_{i \in \Om}$. For all $i \in \Om$ define simple roots and coroots by
\begin{equation} \al_i = x_i - \zeta x_{i+1}, \qquad \al_i^\vee = x_i^* - \zeta^{-1} x_{i+1}^*. \end{equation}
This induces the action
\begin{equation} \label{zetarefrep} s_i(x_i) = \zeta x_{i+1}, \quad s_i(x_{i+1}) = \zeta^{-1} x_i, \quad s_i(x_j) = x_j \text{ if } j \notin \{i,i+1\}. \end{equation}
The Cartan matrix is
\begin{equation} \label{slnzeCartan} \left( \begin{array}{cccccc}
2 & -\zeta & 0 & \cdots & 0 & -\zeta^{-1} \\
-\zeta^{-1} & 2 & -\zeta & \cdots & 0 & 0 \\
0 & -\zeta^{-1} & 2 & \cdots &  & 0 \\
\vdots & \vdots & \vdots & \ddots & -\zeta & 0 \\
0 & 0 &   & -\zeta^{-1} & 2 & -\zeta \\
-\zeta & 0 & 0 & \cdots & -\zeta^{-1} & 2
\end{array} \right), \end{equation}
though in the special case $n=2$ the Cartan matrix is
\begin{equation} \label{sl2zeCartan} \left( \begin{array}{cc}
2 & -(\zeta+\zeta^{-1}) \\
-(\zeta+\zeta^{-1}) & 2 \end{array} \right). \end{equation}
Its determinant is $2 - \zeta^n - \zeta^{-n}$.
\end{example} 
\begin{defn}\label{defn-deformed affine realization}
When $\kk'=\Z$, we shall denote the realization from Example \ref{defn:Vz} by $\Lambda_{\zeta}$ and call it \emph{the (integral) deformed affine realization} of type $A_{n-1}$. 
For any commutative $\Z[\zeta,\zeta^{-1}]$-algebra $A$, we shall refer to the same realization but with $\Bbbk = A$ as the \emph{deformed affine realization over $A$}, and denote it by $\Lambda_A$. Note that $\Lambda_A=\Lambda_{\zeta}\otimes_{\Z[\zeta,\zeta^{-1}]} A.$
\end{defn}


The following lemma is an easy consequence of the definitions. 
\begin{lemma}\label{lemma-deformed affine restricted to Sn}
    The restriction of the (integral) deformed affine realization to $W_{\hh{0}} \cong S_n$ is isomorphic to the permutation realization $\bigoplus_{k=1}^n\Z[\zeta,\zeta^{-1}]y_k$, by sending $y_k \mapsto \zeta^{k-1}x_k$. 
    
    % \BE{why not $y_k = \zeta^k x_k$? Isn't that nicer?}\PT{No specific reason. I just liked my $y_1$ to be $x_1$.}
\end{lemma}



The Dynkin diagram of $W_{\aff}$ admits dihedral symmetry, and these symmetries extend to $\Lambda_{\zeta}$, see \cite[Definition 2.10]{EJY1}.

\begin{defn}We have the \emph{rotation operator} \label{defn-rotation operator}$\sigma$, which is the Dynkin diagram automorphism
$\sigma: \Om \rightarrow \Om: i \mapsto i+1$. 
It acts on $\Lambda_{\zeta}$ by $\sigma(x_i) = x_{\sigma(i)}$ and $\sigma(\zeta) = \zeta$.
We also have the \emph{reflection operator} \label{defn-reflection operator} $\tau: \Om \rightarrow \Om: i \mapsto -i$. It acts on $\Lambda_{\zeta}$ by $\tau(x_i) = x_{1-i}$ and $\tau(\zeta) = \zeta^{-1}$. 
Note that $\tau$ does not preserve the roots; it satisfies $\tau(\alpha_i) = -\zeta^{-1} \alpha_{\tau(i)}$. \end{defn}

Note that $\tau$ fixes the affine vertex, and acts on the finite Weyl group by its traditional Dynkin automorphism (conjugation by the longest element). However, the isomorphism of Lemma \ref{lemma-deformed affine restricted to Sn} does not quite intertwine $\tau$ with the action of the longest element, being off by a multiplicative factor of $\zeta^{n-1}$.





%============================================================
\subsection{Balancedness}
\label{subsec-balanced}
%============================================================
We recall the concept of balancedness.

\begin{defn} (\cite[Definition A.4]{ECathedral}) Let $(W,S)$ be a Coxeter system equipped with a realization, and let $s, t \in S$. It is proven in \cite[\S A.2]{ECathedral} that, whenever $m_{st} = 2k+1$ is odd, one has
\begin{equation} \label{eq:lambdadef} (st)^k(\alpha_s) = \lambda \alpha_t, \quad (ts)^k(\alpha_t) = \lambda^{-1} \alpha_s, \end{equation}
for some scalar $\lambda \in \Bbbk^\times$. The realization is \emph{balanced for} $\{s,t\}$ if $\lambda = 1$. There is a separate definition for $m_{st}$ even that we do not recall here, and it is vacuous when $m_{st} = 2$. We say the realization is \emph{balanced} if it is balanced for all pairs $\{s,t\}$ with $m_{st} < \infty$. \end{defn}

When $m_{st} = 3$, being balanced for $\{s,t\}$ is equivalent to the corresponding entries in the Cartan matrix being $-1$. The deformed affine realization is not balanced.

Under standard notions of what ``positive roots'' should mean, two colinear positive roots should be equal, and both $\alpha_t$ and $(st)^k(\alpha_s)$ should be positive roots. These two requirements are in conflict when the realization is not balanced.

%============================================================
\subsection{Singular Soergel Bimodules from Realizations}
\label{subsec- SSBim}
%============================================================
In \cite{Soer07}, Soergel constructed a category $\SBim$ of bimodules to categorify the Hecke algebra, starting from a nice representation of a Coxeter group $(W,S)$. Williamson, in \cite{WillSingular}, generalized this construction to obtain a 2-category\footnote{In this paper, a 2-category refers to a weak 2-category/bicategory and not necessarily a strict 2-category.} $\SSBim$ of singular Soergel bimodules to categorify the Hecke algebroid\footnote{The Hecke algebroid is sometimes referred to as the Schur algebroid in the literature.}.  We recall these constructions over the next few sections.

 % \PT{Modified $\Lambda^*$ to $\Lambda$ for $R$. Check that the usage is consistent across the paper.}

Given a realization $(\Lambda,\Lambda^*,\{\al_s\},\{\al_s^\vee\})$ over a base (integral domain) $\kk$, we can construct the graded\footnote{All gradings in this paper are over $\Z$.} $\kk$-algebra $R:= \Sym_\kk \Lambda$ (with $\Lambda$ in degree 2). 
The algebra $R$ inherits a $W$-action from $\Lambda$ which preserves the grading. 
For $I\subset S$ we set 
\[ R^I:=R^{W_I}=\{ f \in R: s(f)=f \text{\; for all }s \in I \}.\]

\begin{defn}\label{def sbsbim}
Let $\SBSBim$, the $2$-category of \emph{singular Bott-Samelson bimodules}, be defined as follows. The objects are finitary subsets $I \subset S$, identified with the rings $R^I$. The morphism category $\Hom(I,J)$ will be a full subcategory of graded $(R^J,R^I)$-bimodules, itself denoted $\bModgr{R^J}{R^I}$. The $1$-morphisms in $\SBSBim$ are generated monoidally by the induction bimodule
$\Bimod{R^I}{R^J}{R^I}$ and the shifted restriction bimodule $\Bimod{R^J}{R^I}{R^I} (\ell(J)-\ell(I))$ for the extensions $R^J \subset R^I$ whenever $I \subset J$. \end{defn}

Before continuing we state some grading conventions. For a graded bimodule $M=\bigoplus_{k \in \Z}M_k$ with $k$-th graded component denoted $M_k$, the shifted bimodule $M(j)$ is obtained by modifying the grading so that $M(j)_k=M_{k+j}$.

For $M,N$ being $(R^I, R^J)$-bimodules, $\Hom_{(R^I,R^J)}(M,N)$ consists of bimodule morphisms $M\rightarrow N$ which preserve the grading. 
We also have the notion of a \emph{graded Hom space}. 
\begin{defn}\label{defn- graded Hom}
   The \emph{graded Hom space} $\Hom^\bullet(M,N)$ for $M, N \in \bModgr{R^I}{R^J}$ is defined as 
   $$\Hom^\bullet(M,N):=\bigoplus_{k\in \Z} \Hom_{(R^I,R^J)}(M,N(k)).$$
\end{defn}
For a Laurent polynomial $p=\sum p_iv^i \in \Z_{\geq0}[v,v^{-1}]$ with positive integer coefficients and $M \in \bModgr{R^I}{R^J}$, we define $M^{\oplus p} \in \bModgr{R^I}{R^J}$ to be
$$M^{\oplus p}:=\bigoplus_{i\in \Z}M(-i)^{\oplus p_i}.$$


\begin{defn}
    Let $A$ be a ring and $M$ be a graded $A$-module. Then we say that $M$ is a \emph{graded free} $A$-module if $M\cong A^{\oplus p}$ for a Laurent polynomial $p \in \Z_{\geq 0}[v,v^{-1}]$ with positive integer coefficients, in which case $p$ is called the \emph{graded rank} of $M.$
\end{defn}


Returning to $\SBSBim$, we can describe the 1-morphisms in a more combinatorial manner.

\begin{defn}
Associated to a singular multistep expression as in \eqref{multistepprototype}, we associate the \emph{Bott-Samelson bimodule}
$$\BS(IK_{\bullet}):= R^{I_0}\otimes_{R^{K_1}} \otimes R^{I_1}\otimes \ldots \otimes_{R^{K_d}}R^{I_d}(\sum_{i=1}^d \ell(K_i)-\ell(I_i)) \in \Hom_{\SBSBim}(I_d,I_0).$$
\end{defn}

\begin{example}
    For $S_4$ with simple reflections $\set{s,t,u}$, we have a multistep expression
    $$[[t \subset stu \supset s\subset st]].$$
    The corresponding Bott-Samelson bimodule is the graded $(R^t, R^{st})$-bimodule
    $$R^t \otimes_{R^{stu}}R^s(5).$$
\end{example}


\begin{defn} The $2$-category $\SSBim$ of \emph{singular Soergel bimodules} is the Karoubi envelope (of the graded, additive completion) of $\SBSBim$. For finitary subsets $I, J \subset S,$ we use the notation $\SBimod{I}{J}$ to refer to the category $\Hom_{\SSBim}(J,I)$. 
\end{defn}
Given two graded rings $R_1$ and $R_2$, the category $\bModgr{R_1}{R_2}$ is Karoubi closed. 
Since $\Hom_{\SBSBim}(J,I)$ is a full subcategory of $\bModgr{R^I}{R^J}$, we have that $\SBimod{I}{J}=\Hom_{\SSBim}(J,I)$ is also naturally a full subcategory of $\bModgr{R^I}{R^J}$.

We will see in \cref{thm-categorification} that (under certain assumptions) the indecomposable bimodules in $\SBimod{I}{J}$ are in bijection (up to isomorphism and grading shift) with double cosets $p \in W_I \backslash W / W_J$. Let $p$ correspond to the indecomposable bimodule $B_p$. Whenever $IK_{\bullet}$ is a reduced expression for $p$ as in \cref{def:reducedexpression}, $B_p$ will be a direct summand with multiplicity one inside $\BS(IK_{\bullet})$.

% In particular, for finitary subsets $I, J \subset S$, $\bModgr{R^I}{R^J}$ is Karoubi closed.

% \begin{remark}
%     We can also think of $\SSBim$ as a 2-full sub-2-category of the 2-category of all 1-categories, where we associate to a finitary subset $I \subset S$ the 1-category Mod$^{\text{gr}}(R^I)$ of graded modules over $R^I$, and the generating 1-morphisms are sent to the corresponding induction functors or shifted restriction functors.
% \end{remark}


%=========================================================
\subsection{The Hecke Algebroid}
\label{subsec-Hecke}
%=========================================================

The Hecke algebroid associated to a Coxeter system $(W,S)$ is a relative version of the Hecke algebra associated to all pairs of finitary subsets $I, J \subset S$. One may think of it as a $\Z[v,v^{-1}]$-linear category with objects  indexed by the finitary subsets of $S$. 
We now go into the details of the construction.

Let $\HH$ denote the Hecke algebra associated to $(W,S)$. This is an algebra over $\Z[v,v^{-1}]$. Let $H_w$ denote the standard basis element and $\KL_w$ the Kazhdan-Lusztig basis element corresponding to $w \in W$. The Kazhdan-Lusztig basis is typically complicated, but not when $w = w_J$ for $J$ finitary, where we have
\[ \KL_{w_J} = \sum_{w \in W_J} v^{\ell(w_J) - \ell(w)} H_w. \]
Note that $\KL_{w_J}\KL_{w_J}=\pi(J)\KL_{w_J}$ for all finitary $J \subset S$, where $\pi(J)=\displaystyle \sum_{w\in W_J}v^{\ell(w_J)-2\ell(w)}.$ Thus $\KL_{w_J}$ is a quasi-idempotent in $\HH$. One should think of the Hecke algebroid as something like a (partial) Karoubi envelope associated to these quasi-idempotents.

For all pairs of finitary subsets $I,J \subset S,$ define:
\begin{align*}
    {}^I\HH &= \KL_{w_I}\HH, \quad \text{a right ideal,}\\
    \HH^J &= \HH\KL_{w_J}, \quad \text{a left ideal,}\\
    {^I}\HH^{J} &= {}^I\HH\cap\HH^J.
\end{align*}
Given finitary subsets $I, J, K$, we define a multiplication:
\begin{align*}
    \ast_J :{}^I\HH^J\times {}^J\HH^K \longrightarrow {}^I\HH^K:
    (h_1,h_2) \longmapsto h_1 \ast_J h_2 := \frac{1}{\pi(J)}h_1h_2.
\end{align*}
The operation $\ast_J$ is well-defined as the numerator is divisible by the denominator.
Indeed, for $h_1 \in \Hecke{I}{J}$ and $h_2 \in \Hecke{J}{K}$, we may write $h_1=\KL_{w_I}y_1=x_1\KL_{w_J}$ and $h_2=\KL_{w_J} y_2=x_2\KL_{w_K}$ for some $x_i, y_i \in \HH$, so that 
$$h_1\ast_J h_2=\frac{1}{\pi(J)}h_1h_2=\frac{1}{\pi(J)} x_1\KL_{w_J}\KL_{w_J}y_2=x_1\KL_{w_J}y_2=\KL_{w_I}y_1y_2=x_1x_2\KL_{w_K} \in \Hecke{I}{K}.$$

\begin{defn}\cite[Definition 2.3.1]{WillSingular}
    The Hecke algebroid $\HAbd$ is the following $\Z[v,v^{-1}]$-linear category:
    \begin{itemize}
        \item Objects are given by finitary subsets $I \subset S$;
        \item For finitary subsets $I, J \subset S$, $\Hom(J,I):={}^I\HH^J$, and composition ${}^I\HH^J\times {}^J\HH^K \rightarrow {}^I\HH^K$ is given by $\ast_J$.
    \end{itemize}
\end{defn}
As explained in \cite[section 2.3]{WillSingular}, $\Hecke{I}{J}$ has a standard basis $\{\HAB{I}{J}{p}: p\in W_I\backslash W/W_J\}$ and a Kazhdan-Lusztig (KL) basis $\{\KLB{I}{J}{p}:p \in W_I\backslash W/W_J \}$ where basis elements are indexed by double cosets $p \in W_I\backslash W/ W_J$. As usual, the KL basis is uniquely determined by two properties: self-duality under the bar involution (inherited from the Hecke algebra), and the change of basis formula
\begin{equation} \KLB{I}{J}{p} = \HAB{I}{J}{p} + \sum_{q < p} \textrm{coeff} \cdot \HAB{I}{J}{q} \end{equation}
where each coefficient lives in $v \Z[v]$. Here $<$ indicates the Bruhat order on double cosets.

When the parabolic subgroups $I$ and $J$ are understood, we write $H_p$ instead of $\HAB{I}{J}{p}$, and $\KL_{p}$ instead of $\KLB{I}{J}{p}$.

\begin{remark}
    When $(W,S)$ is the affine Weyl group associated to a simply connected complex semisimple group $G$, and $I$ is the finitary subset such that $W_I$ is the Weyl group associated to $G$, then $\End(I)=({}^I\HH^{I},\ast_I)$ coincides with the spherical Hecke algebra of $G$ (cf. \cref{lemma affine spherical Hecke in Hecke algebroid})
\end{remark}

\begin{remark}
One may also provide the following alternate description for the Hecke algebroid.
For a finitary subset $I\subset S$, let $\HH_I$ denote the Hecke algebra of $(W_I,I)$ thought of as a subalgebra of $\HH$, and let $\triv_I: \HH_I \rightarrow \Z[v,v^{-1}]$ denote the trivial representation (which sends the standard basis elements $H_{w}$ to $v^{-\ell(w)}$, and hence sends $\KL_{w_I}$ to $\pi(I)$). 
Then for finitary $I,J \subset S$, ${}^I\HH^J$ can be identified with $\Hom_{\HH}(\triv_J\otimes_{\HH_J}\HH, \triv_I\otimes_{\HH_I}\HH)$ and $\ast_J:{}^I\HH^J\times {}^J\HH^K \rightarrow {}^I\HH^K$ coincides with composition of (right) $\HH$-module morphisms via this identification.
\end{remark}

When $I \subset J$ one defines ${}^I H^J \in {}^I \HH^J$ to be equal to $\HAB{I}{J}{p}$, associated to the double coset $p$ containing the identity element. As $p$ is minimal in the Bruhat order we have $\HAB{I}{J}{p} = \KLB{I}{J}{p}$, and the underlying element in the Hecke algebra is $\KL_{w_J}$.  Define ${}^J H^I$ similarly. These are the \emph{standard generators}, as Williamson calls them, and they generate the Hecke algebroid. In \cite[Proposition 2.3.3]{WillSingular} one can find a rule for computing the action of standard generators on the standard basis.

Let $IK_{\bullet}$ be a multistep expression as in \eqref{multistepprototype}. We define the element $\bsh(IK_{\bullet}) \in {}^{I_0} \HH^{I_d}$ as the corresponding product of standard generators
\begin{equation} \bsh(IK_{\bullet}) = {}^{I_0} H^{K_1} \star_{K_1} {}^{K_1} H^{I_1} \star_{I_1} {}^{I_1} H^{K_2} \star_{K_2} \cdots \star_{K_d} {}^{K_d} H^{I_d}. \end{equation}
One can use a similar formula to define $\bsh(I_{\bullet})$ for singlestep expressions as in \eqref{singlestepprototype}. Whenever $I_{\bullet}$ is adapted to $IK_{\bullet}$ (see \cref{rmk:adapted}) one has $\bsh(I_{\bullet}) = \bsh(IK_{\bullet})$.

By iterating the rule from \cite[Proposition 2.3.3]{WillSingular}, one obtains an expansion of $\bsh(I_{\bullet})$ in the standard basis, of the form
\begin{equation} \label{eq:expandbs} \bsh([I_0,I_1,\ldots,I_d]) = \sum_{[p_0,p_1,\ldots,p_d]} \textrm{coeff} \cdot H_{p_d}. \end{equation}
The sum ranges over sequences of double cosets $p_i \in W_{I_0} \backslash W / W_{I_i}$ such that either $p_i \subset p_{i+1}$ or $p_i \supset p_{i+1}$, and such that $p_0$ contains the identity element. The coefficient is a certain polynomial in $\N[v,v^{-1}]$. One feature of this expansion is that, when $I_{\bullet}$ is a reduced expression for $p$, then the coefficient of $H_p$ will be $1$. We will not need the details except in the proof of \cref{lem:dim2}.

%=========================================================
\subsection{Categorification results and Hom formula}
\label{subsec-categorification}
%=========================================================
In section \ref{subsec- SSBim}, we defined (following \cite{WillSingular}) a 2-category constructed from a given realization of a Coxeter group $(W,S)$. 
Under nice enough assumptions (as in \cite{WillSingular}), or under a suitable modification of $\SSBim$ as in \cite{Abe} (see \cref{rmk-Abe SSBim}), the so-called Soergel-Williamson categorification theorem and Hom formula both hold. These assumptions do not hold for the deformed affine realization, so we will need to extend the known results. Let us first recall the results in the literature, and discuss assumptions afterwards. Throughout this section, we shall assume that the base $\kk$ is an infinite field.

\begin{theorem}[Categorification Theorem] \cite[Theorem 1, Theorem 2]{WillSingular}, \cite[Theorem 1.1]{Abe} \label{thm-categorification}
    If \cref{willassume} \hyperref[assumption 1]{(1)}, \hyperref[assumption 2]{(2)} and \hyperref[assumption 3]{(3)} hold, we have the following.
\begin{enumerate}
\item There exists a $\Z[v,v^{-1}]$-linear isomorphism ${}_{I}\ch_{J}\colon [\SBimod{I}{J}]\xrightarrow{\sim}{}^I{\mathcal{H}}^J$ such that the diagram
\[
\begin{tikzcd}\relax
[\SBimod{I}{J}]\times [\SBimod{J}{K}]\arrow[r,"\otimes_{J}"]\arrow[d,"{}_{I}{\ch}_{J}\times {}_{J}{\ch}_{K}"] &\relax [\SBimod{I}{K}]\arrow[d,"{}_{I}{\ch}_{K}"]\\
{}^{I}{\mathcal{H}}^{J}\times{}^{J}{\mathcal{H}}^{K}\arrow[r,"*_{J}"] & {}^{I}{\mathcal{H}}^{K}
\end{tikzcd}
\]
is commutative.
\item For each $p\in W_{I}\backslash W/W_{J}$ there exists an indecomposable object
%${}_{I}{B}_{J}(x)\in \SBimod{I}{J}$
$B_p \in \SBimod{I}{J}$
such that
% ${}_{I}{\ch}_{J}({}_{I}{B}_{J}(x))\in {}^{I}{H}^{J}_{x} + \sum_{y < x}\Z[v,v^{-1}]{}^{I}{H}^{J}_{y}$.
${}_{I}{\ch}_{J}(B_p)\in {}^{I}{H}^{J}_{x} + \sum_{y < x}\Z[v,v^{-1}]{}^{I}{H}^{J}_{y}$.
Moreover, it is unique up to isomorphism.
\item For any indecomposable object $B\in \SBimod{I}{J}$ there exists $p\in W_{I}\backslash W/W_{J}$ and $n\in \Z$ such that
% $B\simeq {}_{I}{B}_{J}(x)(n)$.
$B \simeq B_p(n)$.
\item If a singular multistep expression $IK_{\bullet}$ is a reduced expression for the double coset $p \in W_{I}\backslash W/W_{J}$, then $B_p$ is a summand (with multiplicity 1) of the Bott-Samelson bimodule $\BS(IK_{\bullet})$. All other summands are isomorphic to $B_q(n)$ for some $n \in \Z$ and some $q \le p$ in the Bruhat order on double cosets (see \cite[Proposition 3.3]{KELPBruhat}).
\end{enumerate}
\end{theorem}

Whenever $I \subset J$, $[[I \subset J]]$ is a reduced expression for the double coset $p$ containing the identity, and the induction bimodule $\BS(I \subset J) := \BS([[I \subset J]]) = R^I$ is indecomposable, so equals $B_p$. We have ${}_{I}{\ch}_{J}(B_p) = {}^I\KL^J_p = \KL_{w_J}$. Similar statements can be made for $[[J \supset I]]$ and the shifted restriction bimodule, which has character $\KL_{w_J}$ as well. The categorification theorem sends the generators of $\SSBim$ to the standard generators of the Hecke algebroid. Property (1) above then implies that the character of $\BS(I_{\bullet})$ is $\bsh(I_{\bullet})$ for any singular expression $I_{\bullet}$.


\begin{remark} For certain nice realizations (typically in characteristic zero) the \emph{Soergel conjecture} holds, which states that ${}_{I}{\ch}_{J}(B_p) = {}^I\KL^J_p$ for all $p$. Whether this holds is not relevant in this paper.
\end{remark}

Before stating the Williamson-Soergel Hom formula, we recall the standard pairing on the Hecke algebra $\HH$.
Let $\bar{( \cdot)}: \ZZ[v,v^{-1}] \rightarrow \ZZ[v, v^{-1}]$ denote the ring automorphism $v \mapsto v^{-1}$. 
This extends to the \emph{bar involution} or the \emph{Kazhdan-Lusztig involution} 
$$\bar{(\cdot)}: \HH \rightarrow \HH: H_w \mapsto H_{w^{-1}}^{-1}.$$
We also have the \emph{Kazhdan-Lustig anti-involution} $\omega\colon \HH\to \HH$  given by 
$$\omega(\sum_{w\in W}a_{w}H_{w}) := \sum_{w\in W}\overline{a_{w}}H_{w}^{-1}.$$
This is an anti-involution and $\omega(\KL_{w_I}) = \KL_{w_{I}}$.
We also have the \emph{standard trace} $\epsilon: \HH \rightarrow \ZZ[v,v^{-1}]$ which extracts the coefficient of $H_1$ when writing an element of $\HH$ in the standard basis. 
Note that $\epsilon(ab)=\epsilon(ba)$ (see, for example, \cite[Corollary 3.17]{EMTW}).
\begin{defn}\label{defn-standard form}
    The \emph{standard form} $(-,-):\HH \times \HH \rightarrow \ZZ[v,v^{-1}]$ is the sesquilinear form 
    $$(a,b):=\epsilon(\omega(a)b) = \epsilon(b\omega(a)).$$
    This form can be extended to the Hecke algebroid as $(-,-):\Hecke{I}{J} \times \Hecke{I}{J} \rightarrow \ZZ[v,v^{-1}]$, \begin{equation} (a,b):= v^{-\ell(J)}\epsilon(b*_J\omega(a)).\end{equation} Note that $\omega(\Hecke{I}{J}) = \Hecke{J}{I}$, so that $b *_J \omega(a)$ makes sense.\footnote{The factor of $v^{-\ell(J)}$ is not always included in the literature.}
\end{defn}

Just as for the Hecke algebra, the standard form on the Hecke algebroid has the following property, the \emph{asymptotic orthonormality} of the KL basis.

\begin{lemma} \label{lem:ortho} Let $p, q \in W_I \backslash W/W_J$. Then
\begin{equation} (\KLB{I}{J}{p}, \KLB{I}{J}{q}) \in \delta_{pq} + v \Z[v]. \end{equation}
\end{lemma}

Now we state the Hom formula. We recommend the reader see \cite[Theorem 4.16]{EKLP} for a reformulation which is often more useful in practice, one which also makes \cref{lem:ortho} more transparent.

\begin{theorem}[Hom formula]\cite[Theorem 7.4.1]{WillSingular}, \cite[Theorem 2.54]{Abe} \label{thm-hom formula}
If \cref{willassume} \hyperref[assumption 1]{(1)}, \hyperref[assumption 2]{(2)} and \hyperref[assumption 3]{(3)} hold, we have the following. 
Let $M_{1},M_{2}\in \SBimod{I}{J}$.
Then $\Hom_{\SBimod{I}{J}}(M_{1},M_{2})$ is a graded free $R^{I}$-module and we have
\[
\grk_{R^{I}} \Hom^\bullet_{\SBimod{I}{J}}(M_{1},M_{2}) = 
({}_{I}{\ch}_{J}(M_{1}), {}_{I}{\ch}_{J}(M_{2}))
\]
\end{theorem}

\begin{remark} The morphism space is also free as a right $R^J$-module with a different graded rank. \end{remark}

We now give an extended example, as the proof of the following lemma. This lemma will play a role in our proof of the well-definedness of the 2-functor $\GS$.

\begin{lemma} \label{lem:dim2} Let $(W,S)$ be the affine Weyl group using conventions from \cref{sec:combo}. Consider the fundword $\ula = (\varpi_{k-1}, \varpi_1)$ for some $2 \le k \le n$, and let $I_{\bullet} = \GS(\ula,0)$, see \cref{def:GSfundword}. Explicitly, we have
\begin{equation} I_{\bullet} = [\hh{k} \supset \hh{k1} \subset \hh{1} \supset \hh{01} \subset \hh{0}]. \end{equation}
Then $\dim \End^0(\BS(I_{\bullet})) = 2$. \end{lemma}

\begin{proof}
By the Hom formula, we should compute the pairing $(\bsh(I_{\bullet}), \bsh(I_{\bullet}))$, which yields the graded rank of the endomorphism space over $R^{\hh{k}}$. This does not agree with the graded dimension of the endomorphism space, but they have the same leading term (the smallest power of $v$) since $R^{\hh{k}}$ is concentrated in non-negative degree with a one-dimensional degree zero space. Thus if the pairing lives in $\Z[v]$, then $\dim \End^0(\BS(I_{\bullet}))$ agrees with the coefficient of $v^0$.

Our goal is to prove that
\begin{equation} \label{dim2bsdecomp} \bsh(I_{\bullet}) = \KL_{\psi_0(\varpi_1 + \varpi_{k-1})} + \KL_{\psi_0(\varpi_k)}. \end{equation}
By the almost orthonormality of the KL basis \cref{lem:ortho}, we deduce that 
\[ (\bsh(I_{\bullet}), \bsh(I_{\bullet})) \in 2 + v \Z[v], \]
whence we deduce that $\dim \End^0(\BS(I_{\bullet})) = 2$.

First assume $k < n$. We expand $\bsh(I_{\bullet})$ in the standard basis. We consider all sequences $[p_0 \supset p_1 \subset p_2 \supset p_3 \subset p_4]$ as in \eqref{eq:expandbs}. There is a unique choice of $p_0$ and $p_1$ and $p_2$: the underlying set of both $p_0$ and $p_1$ is $W_{\hh{k}}$, and $p_2 = W_{\hh{k}} 1 W_{\hh{1}}$. In fact one has
\begin{equation} \bsh([\hh{k} \supset \hh{k1} \subset \hh{1}]) = H_{p_2} \end{equation}
since this is a reduced expression and $p_2$ is minimal in the Bruhat order.

There are two double cosets in $W_{\hh{k}}\backslash W /W_{\hh{01}}$ contained in $p_2$, denoted $p_3 < p'_3$. They satisfy
\begin{equation} \underline{p}_3 = \underline{p}_2 = 1, \qquad \overline{p}_3' = \overline{p}_2 = w_{\hh{k}} w_{\hh{k1}}^{-1} w_{\hh{1}}. \end{equation} 
Finally, $p_3$ and $p_3'$ are contained respectively in $p_4$ and $p_4'$, with
\begin{equation} p_4 = \psi_0(\varpi_{k}), \qquad p_4' = \psi_0(\varpi_1 + \varpi_{k-1}). \end{equation}

As proven in \cref{thm:GSrex}, $I_{\bullet}$ is a reduced expression for $\psi_0(\varpi_1 + \varpi_{k-1})$, so the coefficient of $H_{p_4'}$ in this expansion will be $1$. Meanwhile, according to \cite[Proposition 2.3.3(1)]{WillSingular}, the coefficient of $H_{p_4}$ will be
\begin{equation} v^{\ell(\overline{p_2}) - \ell(\overline{p_3})} v^{\ell(\underline{p_3}) - \ell(\underline{p_4})} \frac{\pi(\hh{0k})}{\pi(\hh{01k})}, \end{equation}
where $\pi$ is the balanced Poincare polynomial of the parabolic subgroup. This ratio of Poincare polynomials is that for $S_k$ over $S_{k-1}$, and is equal to $[k] = v^{k-1} + v^{k-3} + \ldots + v^{-(k-1)}$. The difference $\ell(\overline{p_2}) - \ell(\overline{p_3})$ is equal to $k-1$, whereas the relevant minimal elements are equal to the identity. Thus we have
\begin{equation} \bsh(I_{\bullet}) = H_{p_4'} + v^{k-1} [k] H_{p_4} = H_{\psi_0(\varpi_1 + \varpi_{k-1})} + (1 + v^2 + \ldots + v^{2(k-1)}) H_{\psi_0(\varpi_k)}. \end{equation}

Now we re-express this in the KL basis. Note that $\bsh(I_{\bullet})$ is self-dual (as is any product of the self-dual generators), and $H_{\psi_0(\varpi_k)} = \KL_{\psi_0(\varpi_k)}$ since the double coset is minimal in the Bruhat order, so it is also self-dual. Thus the difference is self-dual, and the coefficient of $H_{\psi_0(\varpi_k)}$ lives in $v \Z[v]$. By the defining property of the KL basis we must have
\begin{equation} \KL_{\psi_0(\varpi_1 + \varpi_{k-1})} = \bsh(I_{\bullet}) - H_{\psi_0(\varpi_k)} = H_{\psi_0(\varpi_1 + \varpi_{k-1})} + (v^2 + \ldots + v^{2(k-1)}) H_{\psi_0(\varpi_k)}. \end{equation}
In particular, \eqref{dim2bsdecomp} holds as desired.

When $k=n$, replace $\varpi_k$ with the zero weight, so that $p_4 = \psi_0(0)$ is the double coset containing the identity. One should also interpret $\hh{0k}$ as $\hh{0}$ in this special case. Otherwise, the proof above works directly.
\end{proof}

\begin{remark} \label{rmk:dim2vsrepthry} On the other side of the Soergel Satake equivalence, the corresponding statement is
\begin{equation} \label{dim2inreptheory} \dim \End(L_{\varpi_{k-1}} \ot L_{\varpi_1}) = 2. \end{equation}
One typically proves this using Schur's lemma and the decomposition
\begin{equation}\label{1kdecompreptheory} L_{\varpi_{k-1}} \ot L_{\varpi_1} \cong L_{\varpi_{1} + \varpi_{k-1}} \oplus L_{\varpi_k}. \end{equation}
(Replace $\varpi_k$ by the zero weight if $k=n$.) The decategorified version of this decomposition is \eqref{dim2bsdecomp}.
Note however that \eqref{dim2inreptheory} holds even in finite characteristic where \eqref{1kdecompreptheory} might fail! Similarly, the decomposition
\begin{equation} \BS(I_{\bullet}) \cong B_{\psi_0(\varpi_1 + \varpi_{k-1})} \oplus B_{\psi_0(\varpi_k)} \end{equation}
need not hold in finite characteristic. We will not rely on such decompositions in this paper, only on the decategorified statements and their implications for the size of morphism spaces.
\end{remark}


\subsection{The Satake isomorphism in terms of the Hecke algebroid}\label{subsec Satake Hecke algebroid}
In this section, we rephrase the Satake isomorphism (which was categorified into the geometric Satake equivalence) as an isomorphism between the Grothendieck group of $\Rep(SL_n)$ and a version of the spherical Hecke algebra for $W_{\ext}$ constructed within the Hecke algebroid.
We then define $\Rep^{\Om}(SL_n)$ (or rather, its quantum analogue), and reinterpret the the Satake isomorphism as a functor into the Hecke algebroid.

% We begin by stating the Satake isomorphism (the decategorification of the Satake equivalence)\PT{Why explicitly add``the decategorification of the Satake equivalence"? The Satake isomorphism is more elementary than the geometric Satake equivalence.} in terms of the Grothendieck group of $\Rep(\SL_n)$, as is standard in the literature \PT{I don't actually state the original Satake isomorphism here, but only the reformulated version in terms of $\HHH$ instead of the spherical Hecke algebra.}. Then we define $\Rep^{\Omega}(\SL_n)$ and state the analogous result about its decategorification.

We retain our notation from \cref{sec:combo}, so that $\Lambda_{\wt}$ denotes the weight lattice of $\mathfrak{sl}_n$, $\Lambda_{\wt}^+$ denotes the subset of dominant weights, $\rot: \Lambda_{\wt} \rightarrow \Om$ denotes the quotient map with respect to the root lattice, and $\psi_a$ denotes the bijection in \cref{def:psia}.
Let $\Z[\Lambda_{\wt}]$ denote the the group algebra of $\Lambda_{\wt}$.
% The finite Weyl group $W_{\hh{0}}$ acts on $\Lambda_{\wt}$, hence acts $\Z$-linearly on $\Z[\Lambda_{\wt}]$.
It is well known that $[\Rep(\SL_n)]=\Z[\Lambda_{\wt}]^W$, the subring of $W$-invariants.
One one side of the Satake isomorphism is $[\Rep(\SL_n)]\otimes_{\Z} \Z[v^{\pm 1}]=\Z[v^{\pm 1}][\Lambda_{\wt}]^W.$

There is a decomposition $$\Rep(SL_n)=\bigoplus_{k \in \Om}\Rep(SL_n)_{k}$$
where $\Rep(SL_n)_{k}$ is the full subcategory of representations with weights lying entirely in the root coset $k$. 
This decomposition naturally induces an $\Om$-grading
\begin{equation}\label{eq grading on representation ring of sl_n}
   \Z[\Lambda_{\wt}]^W= \bigoplus_{k \in \Om}\Z[\Lambda_{\wt}]^W_k,
\end{equation}
which further induces an $\Om$-grading on $\Z[v^{\pm 1}][\Lambda_{\wt}]^W$.



The other side of the Satake isomorphism would then be the spherical Hecke algebra associated to the extended affine Weyl group $W_{\ext}$ of $\PGL_n$.
We do not discuss the spherical Hecke algebra here, but instead work within the Hecke algebroid associated to $W$.

Recall that we had the rotation operator $\sigma$ which was an automorphism of the Dynkin diagram (see \cref{defn-rotation operator}). 
Then $\sigma$ naturally acts on the Hecke algebroid, sending $\HAB{I}{J}{p}$ to $\HAB{\sigma(I)}{\sigma(J)}{\sigma(p)}$.
% $W$, and this action extends to the Hecke algebra $\HH$, such that
% $\sigma(H_w)=H_{\sigma(w)}$ for $w \in W$.
% Moreover, $\sigma$ naturally acts on all subsets of the set of simple reflections, hence extending further to an action on the Hecke algebroid. Explicitly, $\sigma$ sends $\HAB{I}{J}{p}$ to $\HAB{\sigma(I)}{\sigma(J)}{\sigma(p)}$.

\begin{defn}\label{defn the algebra HHH}
Let 
\begin{equation}\label{eq HHH graded algebra}
    \HHH=\bigoplus_{0\leq k \leq n-1}\Hecke{\hh{k}}{\hh{0}}.
\end{equation} 
$\HHH$ is naturally a right module over $\Hecke{\hh{0}}{\hh{0}}$, but we define a $\Z[v,v^{-1}]$-algebra structure on $\HHH$ as follows.
For $f \in \Hecke{\hh{k}}{\hh{0}}$, $g \in \Hecke{\hh{l}}{\hh{0}}$, define
\begin{equation} \label{eq:badmultyuck}
    f\bullet g:= \sigma^l(f)\ast_{\hh{l}}g \in \Hecke{\hh{k+l}}{\hh{0}}
\end{equation}
and extend it linearly to $\HHH$. 
Then $\HHH$ is an $\Om$-graded algebra (with respect to the decomposition in \cref{eq HHH graded algebra}) and has $\Hecke{0}{0}$ as a subalgebra.
\end{defn}

The algebra $\HHH$ is in fact isomorphic to the spherical Hecke algebra associated to $W_{\ext}$, and hence may be used on the other side of the Satake isomorphism. Then Lusztig's work in \cite{LusztigGS} may be reformulated as follows; see \cref{subsec extended affine to affine} for further details.

\begin{theorem}\label{thm Soergel Satake mini}
    There is an isomorphism of $\Z[v^{\pm 1}]$-algebras
    $$\Z[v^{\pm 1}][\Lambda_{\wt}]^W \xrightarrow{\sim} \HHH$$
    which sends $\ch(L_\lambda)$ to $\KLB{\hh{\rot(\lambda)}}{\hh{0}}{\psi_0(\lambda)}$ for $\lambda \in \Lambda_{\wt}^+$, where $\ch(L_\lambda)$ is the character of the irreducible representation of highest weight $\lambda$.
    This isomorphism respects the $\Om$-grading on these algebras given by \cref{eq grading on representation ring of sl_n}, \cref{eq HHH graded algebra}.
\end{theorem}

The algebra structure on $\HHH$ may appear slightly artificial, hence we prefer to work directly with the Hecke algebroid, and do not discriminate between maximal finitary subsets of $S$.

%We have the following analogue for $\Z[\Lambda_{\wt}]^W$.
\begin{defn}\label{defn KRepom}
    The category $\KRepom$ is defined as follows.
    \begin{itemize}
        \item  The objects of $\KRepom$ are elements of $\Om=\Z/n\Z$.
        \item The morphism space $\Hom(l,k)$ is the $\Z[\Lambda_{\wt}]^W_{k-l}$, spanned by the characters of representations whose weights live in the root coset $k-l \in \Om$. Composition is given by multiplication within $\Z[\Lambda_{\wt}]^W.$
    \end{itemize}
    This is naturally a $\Z$-linear category. 
    We denote by $\KRepom \otimes_{\Z}\Z[v^{\pm 1}]$ the $\Z[v^{\pm 1}]$-linear category with same objects, but with morphism spaces given by $\Hom_{\KRepom}(l,k)\otimes_{\Z}\Z[v^{\pm 1}]=\Z[v^{\pm 1}][\Lambda_{\wt}]^W_k$.
\end{defn}

The following theorem is equivalent \cref{thm Soergel Satake mini}, after translating the unnatural multiplication of \eqref{eq:badmultyuck} into the more natural multiplication in the Hecke algebroid $\HAbd$.

\begin{theorem}\label{thm Soergel Satake}
    There is a functor $\KRepom \rightarrow \HAbd$ which
    sends the object $k \in \Om$ to the maximal finitary subset $\hh{k}$,
    and sends $\ch(L_{\lambda})\in \Hom(l,k)$ to $\KLB{\hh{k}}{\hh{l}}{\psi_l(\lambda)} \in \Hecke{\hh{k}}{\hh{l}}$.
    The induced functor 
    $$\KRepom \otimes_{\Z}\Z[v^{\pm 1}] \rightarrow \HAbd$$ is fully faithful, with essential image the set of maximal finitary subsets of $S$.
\end{theorem}

We now discuss the categorification of $\KRepom$, which appears on one side of the Soergelified Satake equivalence.

In \cite[Section 4.2]{EQuantumI}, it is explained how to construct a 2-category from an \emph{$H$-graded-monoidal category} (\cite[Definition 4.5]{EQuantumI}) for a finite abelian group $H$. 
As seen before, the category $\Rep(SL_n)$ is naturally an $\Om$-graded-monoidal category, to which applying the construction yields a 2-category $\Rep^{\Om}$.
Similarly, for generic $q$, $\Rep(U_q(\mathfrak{sl}_n))$ is also an $\Om$-graded-monoidal category, to which applying the construction yields a 2-category $\qRep$. 
We recall the precise definition of $\qRep$ below, the definition of $\Rep^{\Om}$ is similar.

\begin{definition}\label{defn- qRep Omega}
    The 2-category $\qRep$ is defined as follows.
    \begin{itemize}
        \item The set of objects in the 2-category is $\Om$. 
        \item The Hom-category $\Hom(l,k)$ is the full subcategory %$\Rep(U_q(\mathfrak{sl}_n))_{k-l}$ 
        of $\Rep(U_q(\mathfrak{sl}_n))$ (defined over $\Q(q)$) consisting of representations whose weights live in the root coset $k-l$.
        \item Composition of 1-morphisms is given by tensor product of representations.
    \end{itemize}
\end{definition}

\begin{remark}
    In \cite{EQuantumI}, the set of objects is taken to be the set of ``removable'' vertices of the Dynkin diagram, which is a torsor over $\Om$. This set is then shown to be in canonical bijection with $\Om$, with the affine vertex identifying with $0$ (see \cite[Claim 4.13]{EQuantumI}).
\end{remark}

Recall that at the level of Grothendieck rings, $[\Rep(U_q(\mathfrak{sl}_n))]=[\Rep(SL_n)]=\Z[\Lambda_{\wt}]^W$ as $\Om$-graded rings. The following lemma is then a consequence of \cref{defn KRepom}, \cref{defn- qRep Omega} and this observation.

\begin{lemma}\label{lemma decategorification of Rep-omega}
    There are canonical isomorphisms $[\Rep^{\Om} ]=[\qRep]=\KRepom$.
\end{lemma}

In view of \cref{lemma decategorification of Rep-omega}, \cref{thm Soergel Satake} is the decategorified version of the geometric Satake functor when interpreted in terms of singular Soergel bimodules, or its quantization $\GS_\zeta$ thereof, constructed in \cref{subsec the functor}.

\subsection{Assumptions on a realization}

% Not all of the assumptions \hyperref[assumption 1]{(1)}, \hyperref[assumption 2]{(2)} and \hyperref[assumption 3]{(3)} below will hold in our case, so we also explain how to weaken one of the assumptions so that these results still hold for our deformed affine realization in section \ref{subsec-deformed realization}.

Now we discuss the assumptions made by Williamson in his proofs of these results, and how to weaken those assumptions.
First we recall the notion of reflection faithfulness for a representation $V$ of a Coxeter group $(W,I)$ as introduced in \cite{Soer07} and later used in \cite{WillSingular}.

\begin{defn}\label{defn-reff over field}
    Let $(W,S)$ be a Coxeter system, and $\kk$ be a field. A \emph{reflection faithful} representation of $W$ is a finite
dimensional representation $V$ of $W$ over $\kk$ such that:
\begin{itemize}
\item The representation is faithful;
\item 
We have $\codim V^w = 1$ if and only if $w$ is a reflection.
\end{itemize}
\end{defn}

\begin{lemma}\label{lemma reff implies (3')}
   Let char $\kk \neq 2$ and let $V$ be a reflection faithful representation of $(W,S)$. Then for any reflection $t \in W$, there exists a pair $(h_t,v_t)\in V^* \times V$ such that 
   \[t(\lambda)=\lambda - 2h_t(\lambda)v_t\]
   for all $\lambda\in V$. The pair $(h_t, v_t)$ is unique up to the equivalence relation $(h_t, v_t) \sim (a h_t, a^{-1} v_t)$ for units $a \in \kk^{\times}$.
\end{lemma}

\begin{proof}
When char $\Bbbk$ is not $2$, any involution acts diagonalizably with eigenvalues $1$ and $-1$.
When $V$ is a reflection faithful representation of $W$ and $t$ is a reflection, the second condition in \cref{defn-reff over field} then implies that $V$ has a one dimensional $-1$-eigenspace for $t$. Choose any non-zero element $v_t$ in the $-1$-eigenspace, and let $h_t \in V^*$ be the unique element which vanishes on $V^t$ and takes the value $1$ on $v_t$. By verifying the statement for $V^t$ and for $v_t$, one concludes that $t(\lambda)=\lambda - 2h_t(\lambda)v_t$. The uniqueness statement is straightforward.
\end{proof}

Of course, one should think of $(h_t, v_t)$ as a root and (half a) coroot attached to $t$.





% It follows that there exists $v_t \in V$ (which may be chosen to be any non-zero element of the $-1$-eigenspace) and $h_t \in V^*$ (which vanishes on $V^t$ and takes the value $2$ \BE{you said the value 1 before)} on $v_t$) such that the action of $t$ is given by $t(\lambda)=\lambda - h_t(\lambda)v_t$.
% Note that any such $h_t$ (resp. $v_t$) has to be a $-1$-eigenvector for $t$ in $V^*$ (resp. $V)$, and the pair $(h_t,v_t)\in V^* \times V$ is necessarily unique up to rescaling $v_t$ by a unit in $\kk^\times$ and rescaling $h_t$ by its inverse.
% We then have the following lemma:

Let $T$ denote the set of reflections in $W$. For $t,t'\in T$, fix some $h_t, h_{t'}$ as in the above lemma. If $h_t, h_{t'}$ are not linearly independent, then $V^t=\ker(h_t)=\ker(h_{t'})=V^{t'}$, so that $\codim V^{tt'}\leq 1$. 
Since $tt'$ is not a reflection, reflection faithfulness of $V$ implies that $V^{tt'}=V$, and we have that $t=t'$ by faithfulness of the representation $V$.
Hence we have the following lemma (\cite[Lemma 4.1.4]{WillSingular}):
\begin{lemma}\label{lemma coroots are linearly independent}
     Let char $\kk \neq 2$ and $V$ be a reflection faithful representation of $(W,S)$. The elements of $\set{h_t: t \in T} \subset V^*$ as above are pairwise linearly independent.
\end{lemma}

Theorems \ref{thm-categorification} and \ref{thm-hom formula} as proven in \cite{WillSingular} assume that the following properties hold.

\begin{property} \label{willassume} The following are properties of a finite-dimensional realization $\Lambda$ of $(W,S)$ over a field $\kk$. As usual, we regard $\Lambda^*$ as a $W$-module via the contragredient action, and let $R = \Sym_{\Bbbk}(\Lambda)$.

\begin{enumerate}
    \item $\Lambda^*$ is a reflection faithful representation of $(W,S)$, and  $\kk$ is an infinite field of characteristic $\ne 2$. \label{assumption 1}
    \item For all finitary $I \subset S$, $R$ is graded free over $R^I$, and one has an isomorphism of graded $R^I$-modules: 
    $$R \cong \Tilde{\pi}(I)R^I,$$
    where $\tilde{\pi}(I):= \displaystyle \sum_{w \in W_I}v^{-2\ell(w)}.$\label{assumption 2}
    \item One may choose a family of pairs $\{(h_t, v_t) \colon t \in T\}$ as in \cref{lemma reff implies (3')}, such that
    $$xh_s = h_t \text{ if } xsx^{-1} = t.$$
    \label{assumption 3}
\end{enumerate}
\end{property}

\begin{remark} The statement of \cref{willassume} \hyperref[assumption 1]{(1)} was phrased in terms of $\Lambda^*$ rather than $\Lambda$ to be consistent with Williamson's conventions. Regardless, when comparing with Lie theory, the roots (not the coroots) should be elements of $R$. It is not hard to verify that \hyperref[assumption 1]{(1)} holds for $\Lambda^*$ if and only if it holds for $\Lambda$. The same can be said for \hyperref[assumption 3]{(3)}, but we are not sure whether the same can be said for \hyperref[assumption 2]{(2)}. \end{remark}

\begin{remark} One can think of \cref{willassume} \hyperref[assumption 3]{(3)} as the statement that the realization has something like a ``non-integral root system.'' In \cite[Section 4.1]{WillSingular}, both \cref{willassume} \hyperref[assumption 1]{(1)} and \hyperref[assumption 2]{(2)} are stated as explicit assumptions, but \hyperref[assumption 3]{(3)} is merely asserted as a property which holds. This assertion is incorrect, so we view this property as a third assumption made by Williamson. \end{remark}

\begin{lemma} Property \hyperref[assumption 3]{(3)} fails for the deformed affine realization. \end{lemma}

\begin{proof} Suppose that \hyperref[assumption 3]{(3)} holds, and choose a family of elements $(h_t, v_t)$ for each reflection. Note that $h_{s_1}$ must be in the span of $x_1 - \zeta x_2$; by rescaling the entire family we can assume $h_{s_1} = x_1 - \zeta x_2$. Conjugating $s_1$ by $s_1 s_2$ yields, $s_2$, so we apply $s_1 s_2$ to $h_{s_1}$ and deduce that $h_{s_2} = \zeta x_1 - \zeta^2 x_2$. Continuing to apply $s_i s_{i+1}$ to $h_{s_i}$, we deduce that $h_{s_{n+1}} = \zeta^{n} x_1 - \zeta^{n+1} x_2$. But $s_{n+1} = s_1$, a contradiction. \end{proof}

Let us state a weaker property which a realization might possess. 

\begin{property} We continue the setup of \cref{willassume}.
\begin{enumerate}
    \item[(3')]\label{assumption 3'} 
    One may choose a family of pairs $\{(h_t, v_t) \colon t \in T\}$ as in \cref{lemma reff implies (3')}, such that
    $$xh_s \in \kk^\times \cdot h_t \text{ if } xsx^{-1} = t$$
    where $\kk^\times$ denotes the group of units in $\kk$.  
\end{enumerate}
\end{property}

\begin{lemma}  Property \hyperref[assumption 1]{(1)} implies Property \hyperref[assumption 3']{(3')}. \end{lemma}

\begin{proof} In fact, one can choose any pairs $(h_t, v_t)$ for each reflection $t$, and Property \hyperref[assumption 3']{(3')} will hold. By \cref{lemma reff implies (3')}, one need only show that $xh_s$ is a $-1$-eigenvector for $t$, which is obvious. \end{proof}

We shall see in section \ref{subsec-deformed realization} that our deformed affine realization satisfies assumptions \hyperref[assumption 1]{(1)}, \hyperref[assumption 2]{(2)}, and thus also \hyperref[assumption 3']{(3')}, when working over $\kk=\Q(\zeta)$ for an indeterminate $\zeta.$ 

Now we explain why $\hyperref[assumption 3']{(3')}$ is sufficient for the arguments in \cite{WillSingular} to go through.  
Property $\hyperref[assumption 3]{(3)}$ appears as eq $(4.1.3)$ in \cite{WillSingular} and is used only in the proofs of \cite[Definition/Proposition 4.3.2, Definition/Proposition 5.0.1, and Lemma 7.3.4]{WillSingular}. We recall these briefly below and explain exactly why $\hyperref[assumption 3']{(3')}$ is sufficient. Note that Williamson fixes the family $\{(h_t, v_t)\}$ satisfying $\hyperref[assumption 3]{(3)}$ once and for all; we too fix a family satisfying $\hyperref[assumption 3']{(3')}$.

\begin{prop*}\cite[Definition/Proposition 4.3.2]{WillSingular} \label{deprop-Wil 4.3.2}
Let $X \subset W$ be a finite subset. Consider the subspace
\begin{equation*}
R(X) = \left \{ f = (f_x) \in \bigoplus_{x \in X} R \; \middle | \; 
\begin{array}{c} f_x - f_{tx} \in (h_t) \\ \text{ for all } t \in T
\text{ and } x, tx \in X\end{array}  \right \} \subset \bigoplus_{x \in X} R.
\end{equation*}
Then $R(X)$ is a graded $k$-algebra under componentwise multiplication. 
It becomes an object of $\bMod{R}{R}$ if we define left and right actions of $r \in R$ via
\begin{align*}
 (rf)_x & = r \cdot f_x \\
 (fr)_x & = f_x\cdot x(r)
\end{align*}
for $f = (f_x) \in R(X)$ (where $\cdot$ on the RHS above is multiplication in $R$).
If a pair of subgroups $W_1, W_2 \subset W$ satisfy $W_1X = X = XW_2$ then $R(X)$ carries commuting left $W_1$- and right $W_2$-actions if we define
\begin{align*}
  (uf)_x &= u(f_{u^{-1}x}) & \text{for $u \in W_1$},\\
  (fv)_x &= f_{xv^{-1}} & \text{for $v \in W_2$}.
\end{align*}
\end{prop*}
\begin{remark}
    \cite[Proposition 4.3.2]{WillSingular} as stated also includes that if $X = \{ x \}$ is a singleton then $R(X) \cong R_x$, where $R_x$ denotes the standard bimodule as in \cite[Definition 4.2.1]{WillSingular}, and that if $X = \{x, y \}$ consists of two elements we write $R_{x,y}$ instead of $R(X)$. The former statement is a straightforward consequence of the definitions, and the latter is just notation.
\end{remark}
The proof of this proposition in \cite{WillSingular} only uses $\hyperref[assumption 3]{(3)}$ near the end of the proof to show that the $R(X)$ is stable under the left $W_1$-action. More explicitly, it is shown that if $x, tx \in X$ one has,
\begin{equation*}
  (wf)_x - (wf)_{tx} = w(f_{w^{-1}x} - f_{w^{-1}tx}) \in (w(h_{w^{-1}tw})) = (h_t)
\end{equation*}
where the last equality follows from $\hyperref[assumption 3]{(3)}.$ But the equality of the ideals $(w(h_{w^{-1}tw}))$ and $(h_t)$ also hold with the weaker assumption $\hyperref[assumption 3']{(3')}$. Thus the proposition holds under the assumption $\hyperref[assumption 3']{(3')}$ instead of $\hyperref[assumption 3]{(3)}.$
\vspace{5pt}

Now we proceed to discuss why \cite[Definition/Proposition 5.0.1]{WillSingular} still holds with $\hyperref[assumption 3']{(3')}$. We will continue using the setup in \cite[Definition/Proposition 4.3.2]{WillSingular} discussed above.
\begin{prop*}\cite[Definition/Proposition 5.0.1]{WillSingular}
Let $X, W_1, W_2 \subset W$ be as above.
%Let $I, J \subset S$ be finitary and $p \in \W{I}{J}$.

\begin{enumerate}
\item For all reflections $t \in W_1$ there exists an operator
$f \mapsto \partial_tf$ on $R(X)$, the \emph{left Demazure operator to $t$}, uniquely determined by
\begin{equation*}
  f - tf = 2h_t(\partial_tf) \quad \text{for all $f \in R(X)$.}
\end{equation*}
This is a morphism in $\bMod{R^t}{R}$.

\item For all reflections $t \in W_2$ there exists an operator $f \mapsto f\partial_t$ on $R(X)$, the \emph{right Demazure operator to $t$}, uniquely determined by
\begin{equation*}
  f - ft = (f \partial_t)2h_t \quad \text{for all $f \in R(X)$.}
\end{equation*}
This is a morphism in $\bMod{R}{R^t}$.
\end{enumerate}
\end{prop*}
\begin{proof}
In \cite{WillSingular}, the left Demazure operator $f \mapsto \dd_tf$ is defined by
\begin{equation*}
    (\partial_tf)_x := \frac{f_x - tf_{tx}}{2h_t}, \qquad x \in X
  \end{equation*} 
and the right Demazure operator $f\mapsto f\dd_t$ is defined by
\begin{equation*}
  (f\partial_t)_x := \frac{f_x - f_{xt}}{2x(h_t)},\qquad x\in X.
\end{equation*}
The proof of part (1) of the proposition in \cite{WillSingular} (namely, that the left Demazure operator on $R(X)$ is well-defined, unique, and is a morphism in $\bMod{R^t}{R}$) is completely independent of the assumption $\hyperref[assumption 3]{(3)}$, so part (1) holds.
The explicit check that the right Demazure operator is well-defined as an operator on $R(X)$ is left to the reader in \cite{WillSingular}, so we explain the details of this check and show that it suffices to assume $\hyperref[assumption 3']{(3')}$ instead of $\hyperref[assumption 3]{(3)}$.  Unicity follows from the argument in \cite{WillSingular}.

Let us fix $f \in R(X)$ and $t\in W_2$. To see that $(f\dd_t)_x \in R$ (and not just in $R[x(h_t)^{-1}]$) for $x\in X$, we use that $f_x-f_{xt}=f_x-f_{xtx^{-1}x} \in (h_{xtx^{-1}})=(x(h_t))$. Note that this argument still holds with assumption $\hyperref[assumption 3']{(3')}$ instead of $\hyperref[assumption 3]{(3)}$, since all we need is the equality of the ideals $(h_{xtx^{-1}})$ and $(x(h_t))$. Now we are left to show that $f\dd_t$ in fact lives in $R(X)$ and not just in $\bigoplus_{x\in X}R$. 
In other words, we need to show for any reflection $t' \in W$ and $x \in X$ such that $t'x \in X$,
\begin{equation*}
    (f\dd_t)_x - (f\dd_t)_{t'x} \in (h_{t'}).
\end{equation*}
Adding and subtracting $t'((f\dd_t)_{t'x})$, we see that we need to show
\begin{equation}\label{eq 1 prop williamson demazure operators}
    (f\dd_t)_x-t'((f\dd_t)_{t'x}) - ((f\dd_t)_{t'x}-t'((f\dd_t)_{t'x})) \in (h_{t'}).
\end{equation}
Note that for any $g \in R$ which is $t'$-anti-invariant, $g$ vanishes on $V^{t'}$, so that $g \in (h_{t'})$ (since $(h_{t'})$ is the ideal of functions vanishing on the hyperplane $V^{t'} \subset V$).
In particular, $(f\dd_t)_{t'x}-t'((f\dd_t)_{t'x}) \in (h_{t'})$, so showing \cref{eq 1 prop williamson demazure operators} is equivalent to showing
\begin{equation}\label{eq 2 prop williamson demazure operators}
    (f\dd_t)_x-t'((f\dd_t)_{t'x}) \in (h_{t'}).
\end{equation}
Expanding using definitions and the fact that $t'$ is a reflection, and in particular an involution, we have 
\begin{align}
    (f\dd_t)_x-t'((f\dd_t)_{t'x}) &= \frac{f_x - f_{xt}}{2x(h_t)}-t'\left(\frac{f_{t'x} - f_{t'xt}}{2t'x(h_t)} \right) \nonumber \\
    &=\frac{f_x-t'(f_{t'x}) - (f_{xt}-t'(f_{t'xt}))}{2x(h_t)}. \label{eq-right demazure}
\end{align}

Let $X'=\set{x,xt,t'x,t'xt}$. 
Then $\set{1,t'}X'=X'=X'\set{1,t}$, and we have $\tilde{f}\in R(X')$ whose components are given by those of $f$.
We may then write the numerator above using the left Demazure operators (from part (1) of the proposition, for $\tilde{f} \in R(X')$) as
$$f_x-t'(f_{t'x}) - (f_{xt}-t'(f_{t'xt}))=2h_{t'}\left((\dd_{t'}\tilde{f})_x-(\dd_{t'}\tilde{f})_{xt}\right),$$
so that the numerator is in $(h_{t'}).$
So if $x(h_t)$ is coprime to $h_{t'}$, we have that the RHS of \cref{eq-right demazure} is in $(h_{t'})$ as desired.\\
Now suppose $x(h_t)$ is not coprime to $h_{t'}$. 
From assumption $\hyperref[assumption 3']{(3')}$, we have that $x(h_t) \in \kk^\times h_{xtx^{-1}}$, and reflection faithfulness implies that $\kk h_{xtx^{-1}}=\kk h_{t'}$ if and only if $xtx^{-1}=t'$ (see \cref{lemma coroots are linearly independent}).
Hence we have that $xtx^{-1}=t'$ in this case,
and $x(h_t)=\lambda h_{t'}$ for some $\lambda \in \kk^\times$. 
Then the RHS of \eqref{eq-right demazure} simplifies to $\lambda^{-1}\left((\dd_{t'}\tilde{f})_x-(\dd_{t'}\tilde{f})_{xt=t'x} \right)$ which is in $(h_{t'})$ since we already know that $\dd_{t'}\tilde{f} \in R(X')$ by part (1) of the proposition. 

The statement that $f \mapsto f\dd_t$ is a morphism in $\bMod{R}{R^t}$ follows from the observation that for $a\in R$ and $b \in R^t$,
\begin{align*}
(a(f\dd_t)b)_x=a\left(\frac{f_x - f_{xt}}{2x(h_t)} \right)x(b)= \frac{af_xx(b) - af_{xt}xt(b)}{2x(h_t)}=\frac{(afb)_x - (afb)_{xt}}{2x(h_t)}=((afb)\dd_t)_x.&\qedhere
\end{align*}
\end{proof}

Finally, we recall \cite[Lemma 7.3.4]{WillSingular} and show that it holds under assumption $\hyperref[assumption 3']{(3')}$ instead of $\hyperref[assumption 3]{(3)}$.
In this Lemma, $T$ is the set of reflections in $W$ and $p_-$ is the minimal coset representative of a $(W_I,W_J)$-coset $p$.
\begin{lemma*}\cite[Lemma 7.3.4]{WillSingular}
    The element 
    $$m_p=\prod_{t\in T: tp_-<p_-}h_t$$
    lies in $R^{I\cap p_-Jp_-^{-1}}$.
\end{lemma*}
\begin{proof}
    In the proof of the Lemma in \cite{WillSingular}, it is shown that for $s \in I \cap p_-Jp_-^{-1}$ and $t \in T$ such that $tp_-<p_-$, it follows that $stsp_-<p_-$. Note that this argument is purely within the Coxeter group $W$, and is independent of the representation $V$. With assumption $\hyperref[assumption 3]{(3)}$, it follows that $s(m_p)=m_p$ for all $s\in I\cap p_-Jp_-^{-1}$ since we can think of $m_p$ as a product of the $s$-invariant elements $h_th_{sts}$. We claim that $h_th_{sts}$ is still $s$-invariant with assumption \hyperref[assumption 3']{(3')}. Indeed, by \hyperref[assumption 3']{(3')}, we have that $s(h_{sts})=\lambda h_{t}$ for some $\lambda \in \kk^\times$. Then $h_{sts}=\lambda s(h_t)$, so that $s(h_t)=\lambda^{-1}h_{sts}.$ 
    It follows that 
    \begin{align*}s(h_th_{sts})=s(h_t)s(h_{sts})=\lambda^{-1}h_{sts}\lambda h_{t}=h_th_{sts}. &\qedhere 
    \end{align*}
\end{proof}

Hence we conclude that Theorems \ref{thm-categorification} and \ref{thm-hom formula} hold without requiring \hyperref[assumption 3]{(3)}.

\begin{theorem} Theorems \ref{thm-categorification} and \ref{thm-hom formula} hold when one only requires Properties \hyperref[assumption 1]{(1)} and \hyperref[assumption 2]{(2)} but not \hyperref[assumption 3]{(3)}. \end{theorem}



% 


%=========================================================
\subsection{The deformed affine Type A realization}
\label{subsec-deformed realization}
%=========================================================
In this section, we discuss the deformed affine Type A realization (\cref{defn-deformed affine realization}) in more detail, and prove that Properties $\hyperref[assumption 1]{(1)}$ and $\hyperref[assumption 2]{(2)}$ from section \ref{subsec-categorification} hold when the ground field is $\Q(\zeta)$, so that we can use the categorification theorem and Hom formula from section \ref{subsec-categorification}.

In preparation for future work, we aim to establish this result in reasonable generality. For this purpose it will be helpful to think of realizations over integral domains that are not necessarily fields. Since both \cite{WillSingular}, \cite{Soer07} define ``reflection faithfulness" only over fields (see Definition \ref{defn-reff over field}), we give a generalized definition. There are many generalizations which will specialize to the definition in \cite{WillSingular, Soer07}, but this one is sufficient for our purposes.
\begin{defn}\label{defn-reff}
    $V$ is a \emph{reflection faithful} representation of $(W,S)$ over an integral domain $A$ if:
    \begin{enumerate}
        \item [(a)] $V$ is a faithful representation,
        \item [(b)] $V^w$ is free of finite rank over $A$ for all $w\in W$,
        \item [(c)] $V/V^w$ is free of finite rank over $A$ for all $w \in W$ (i.e., $V^w$ is a direct summand of $V$),
        \item [(d)] rk$_A(V/V^w)=1$ if and only if $w$ is a reflection in $W$.
    \end{enumerate}
\end{defn}


\begin{lemma}\label{reff flat base change}
    Reflection faithfulness is stable under flat base change, i.e., if $A\rightarrow B$ is a flat morphism of commutative domains and $V$ a reflection faithful representation of $(W,S)$ over $A$, then $V_B :=V\otimes_A B$ is a reflection faithful representation over $B$.
\end{lemma}
\begin{proof}
    Observe that for $w\in W$,
    \begin{align*}
        V^w &= \ker(V \xrightarrow{w-1}V), \\
        V/V^w &\cong\Ima(V\xrightarrow{w-1}V). \end{align*}
    Since $A \rightarrow B$ is flat, we have that for any map $\phi$, there are canonical isomorphisms $\ker(\phi)\otimes_A B \cong ker(\phi \otimes_A 1_B)$, $\Ima(\phi)\otimes_A B \cong \Ima(\phi \otimes_A 1_B)$. 
    We then have that
    \begin{align*}
        (V_B)^w &= \ker(V\otimes_A B \xrightarrow{(w-1)\otimes 1_B} V\otimes_A B) = V^w \otimes_A B,\\
        V_B/(V_B)^w &= \Ima(V\otimes_A B \xrightarrow{(w-1)\otimes 1_B} V\otimes_A B)=(V/V^w)\otimes_A B
    \end{align*}
    and (a), (b), (c) and (d) in Definition \ref{defn-reff} follow.
\end{proof}

The following lemma helps us compute fixed point sets. The product of a diagonal matrix and a permutation matrix is called a \emph{rescaled permutation matrix}.

\begin{lemma}\label{lemma rescaled permutation matrix fixed points} Let $\{x_1, \ldots, x_n\}$ be a basis of a free module $V$ over a commutative domain $\Bbbk$, and let $\rho \in GL(V)$ be a rescaled permutation matrix. Then both  the  fixed point set $V^{\rho}$ and the quotient $V/V^{\rho}$ are free. The rank of $V^{\rho}$ is the number of monodromy-free cycles of $\rho$ (as defined in the proof). \end{lemma}

\begin{proof} Let $\rho = w t$ where $w$ is a permutation matrix and $t$ a diagonal matrix. Let $\{1, \ldots, n\} = \coprod \mathcal{O}_k$ be the decomposition into cycles for $w$. As a module over $\rho$, $V = \bigoplus \Span \{x_i\}_{i \in \mathcal{O}_k}$. Each of the statements in the lemma can be proven individually for each cycle, which together proves the statement for the direct sum.
So, by restriction to a given cycle and by renaming our basis, we can assume that $w = (12 \cdots d)$ is a standard $d$-cycle, and that $t$ is the diagonal matrix with entries $(t_1, t_2, \ldots, t_d)$. 
Note that $d=1$ is allowed. Let $c = \prod_{i=1}^d t_i$. We call the cycle \emph{monodromy-free} if $c=1$. It is easy to verify that $\rho^d = c \cdot I$, which has no nonzero fixed points unless $c = 1$. Thus the fixed point set of $\rho$ is zero unless $c=1$.

Assume $c=1$. Then $\rho(x_1) = t_1 x_2$ and $\rho(x_2) = t_1 t_2 x_3$, etcetera. 
We claim that the vector
\[ \sum_{i=1}^d x_i(\prod_{j<i} t_j)\]
is a basis for the (rank one) $\rho$-fixed subspace, and that $\{x_2, \ldots, x_d\}$ descends to a basis for the quotient. We leave the straightforward argument to the reader. 
% \PT{We can even skip this last bit and leave it for the reader:} Indeed, if an element $\sum_{i}c_ix_i \in V^\rho$ for $c_i \in \kk$, we have that 
% $$\sum_i c_ix_i=\rho(\sum_i c_ix_i)=\sum_i c_i t_ix_{i+1}$$
% so that $c_{i+1}=c_it_i$ (for all $i$), and $\sum_ic_ix_i=c_1(\sum x_i(\displaystyle\prod_{j<i}t_j))$. 
% This implies that $V^\rho$ is free, spanned by $v$.
% The statement for $V/V^\rho$ follows.
\end{proof}

Now we consider the deformed affine realization. For simplicity we fix $n \geq 3.$
We briefly recall the basics. 
Let $\Om =\Z/n\Z$, and $W = \tilde{S}_n$, with $S = S_{\aff} = \{s_i\}_{i \in \Om}$. Then
$\Lambda_{\zeta}$ is the realization of $(W,S)$ over $\Z[\zeta, \zeta^{-1}]$ with basis $\{x_i\}_{i \in \Om}$ with $\alpha_i=x_i - \zeta x_{i+1}$ and $\alpha_i^\vee= x_i^*- \zeta^{-1}x_{i+1}^*$, where $\{x_{i}^*\}$ is dual to $\{ x_i\}$. 
The action of $W$ was given by
$$ s_i(x_i)=\zeta x_{i+1}, \hspace{15pt} s_i(x_{i+1})=\zeta^{-1}x_i, \hspace{15pt} s_i(x_j)=x_j \;\text{  otherwise.}$$


% Now we show that for $\kk= \mathbb{Q}(\zeta)$, the representation $V := \Lambda_{\zeta} \otimes_{\Z[\zeta,\zeta^{-1}]}\Q(\zeta)$ of $\tilde{S}_n$ is reflection faithful. 

Let us recall some well-known facts as a way to set notation and context.

\begin{lemma}\label{lemma reflections}
        We have an isomorphism $\Tilde{S}_n \xrightarrow{\sim} S_n \ltimes \Lambda_{\rt}$, realizing the affine Weyl group as affine transformations on the root lattice (of $\mathfrak{sl}_n$), where 
        \begin{align}
        &s_i \mapsto s_i \in S_n \text{ for } i\not=0 \\
        &s_0 \mapsto s_{\alpha_{long}}t_{-\alpha_{long}} = t_{\alpha_{long}} s_{\alpha_{long}}.
        \end{align}
When $\alpha$ is a root of $\mathfrak{sl}_n$, we let $s_{\alpha}$ denote the corresponding reflection in $S_n$ and $t_{\alpha}$ the corresponding translation (an element of $\Lambda_{\rt}$), and we let $t_{k\alpha} = t_{\alpha}^k$. Above, $\alpha_{long}$ is the longest root. Via this isomorphism, the reflections in $W$ get identified with
$\{s_{\alpha}t_{k\alpha}\}$ ranging over $k \in \Z$ and over all positive roots of $\mathfrak{sl}_n$.
\end{lemma}
    
\begin{proof}[Proof] 
The identification of $\tilde{S}_n$ with $S_n\ltimes\Lambda_{\rt}$ is standard and well-known; for example see  \cite[Chapter 4]{HumphreysCoxeter}. The classification of reflections is also well-known but not found in \cite{HumphreysCoxeter}; we sketch the proof.

For $\lambda \in \Lambda_{\rt}$, we shall use the notation $t_{\lambda}$ exclusively  to denote the corresponding element in $W$. Note that $S_n$ normalizes the subgroup $\{t_{\lambda}: \lambda \in \Lambda_{\rt}\}$. 
Explicitly, we have that for $w\in S_n$, $\lambda \in \Lambda_{\rt}$, $wt_{\lambda}=t_{w(\lambda)}w$ where $w(\lambda)$ refers to the usual action of $S_n$ on the root lattice. 
This will be used multiple times without explicit mention for the rest of this section.
    
Writing $w=w't_{\lambda}$ for $w' \in S_n$ and $\lambda \in \Lambda_{\rt}$, we compute that
       \begin{equation}
            ws_{\alpha}t_{k\alpha}w^{-1} = s_{w'(\alpha)}t_{(k-\langle \alpha^\vee, \lambda \rangle)w'(\alpha)}.
        \end{equation}

  % observe that the set of elements of the form $s_{\alpha}t_{k\alpha}$ is stable under conjugation, since
        % \begin{align*}
        %     ws_{\alpha}t_{k\alpha}w^{-1} &= w't_{\lambda}s_{\alpha}t_{k\alpha}t_{-\lambda}w'^{-1}\\
        %     &= w's_{\alpha}t_{s_{\alpha}(\lambda)}t_{k\alpha - \lambda}w'^{-1}\\
        %     &= w's_{\alpha}w'^{-1}t_{w'(k\alpha - \langle \alpha^\vee, \lambda \rangle\alpha)}\\
        %     &= s_{w'(\alpha)}t_{(k-\langle \alpha^\vee, \lambda \rangle)w'(\alpha)}.
        % \end{align*}
 In particular, any reflection $ws_iw^{-1}$ (for $w\in W,\; i \in \ZZ/n\ZZ$) is of the form $s_{\alpha}t_{k\al}$ for some $k \in \ZZ$ and some root $\al$ of $S_n$.
 Hence all reflections in $W$ are indeed of the desired form.
 Note that all the pairings here are for the geometric realization of $S_n$ coming from the root lattice of $\mathfrak{sl}_n$.  


Suppose $n \ge 3$. Conversely, given a root $\al$ of $S_n$, we can find $w' \in S_n$ such that $w'(\al) = \al_1$ whence $w' s_{\al} (w')^{-1} = s_1$. Now we use the fact that $n \ge 3$, so that $s_1(\al_2) = \al_1 + \al_2$, and compute (for any $k \in \Z$) that
\begin{equation}  (t_{-k \al_2} w')(s_{\al} t_{k \al}) (t_{-k \al_2} w')^{-1} = t_{-k\alpha_2}s_1t_{k\alpha_1}t_{k\alpha_2} = s_1. \end{equation}
Thus any element of the form $s_{\al} t_{k \al}$ is a reflection.

The proof is different when $n=2$ since the reflections are not all conjugate, and we leave it to the reader. \end{proof}

% Now, to show the converse, i.e., that any element of the form $s_{\alpha}t_{k\alpha}$ for a root $\alpha$, for $k \in \Z$ is indeed a reflection in $W$, we use the fact that any root of $S_n$ (in the geometric realization) can be acted by an element of $S_n$ to get any other root. 
% So given $s_{\alpha}t_{k\alpha}$, we can choose $w \in S_n$ so that $w(\alpha)=\alpha_1$ and $ws_{\alpha}w^{-1}=s_1$. 
%     But then we have that 
%     \begin{align*}
%         ws_{\alpha}t_{k\alpha}w^{-1}=ws_{\alpha}w^{-1}t_{kw(\alpha)}= s_1t_{k\alpha_1},
%     \end{align*}
%     so it suffices to show that $s_1t_{k\alpha_1}$ is a reflection in $W$. 
%     %The above argument for $\alpha=\alpha_{long}, \;k=1$ shows that %$s_1t_{\alpha_1}$ is a reflection in $W$ (conjugate to $s_0$). 
%     Now, we use the fact that $n\geq 3$, so that $s_2\in S_n$, with $s_1(\alpha_2)=\alpha_1+\alpha_2$. 
%     But then, we observe that
%     $$t_{-k\alpha_2}s_1t_{k\alpha_1}t_{k\alpha_2}=s_1t_{-ks_1(\alpha_2)}t_{k\alpha_1+k\alpha_2}=s_1t_{k(\al_2-s_1(\alpha_2))+k\al_1}=s_1t_0=s_1,$$
%     to conclude that all elements of the form $s_1t_{k\alpha_1}$ are conjugate to $s_1$ in $W$, and hence are indeed reflections.
% \end{proof}

\begin{lemma}\label{lemma-action matrix}
    Identify $GL_{\Z[\zeta, \zeta^{-1}]}(\Lambda_{\zeta})$ with $GL(n,\Z[\zeta, \zeta^{-1}])$ by fixing the basis $\{x_i\}$, and let $\rho: W \rightarrow GL(n,\Z[\zeta, \zeta^{-1}])$ be the map obtained from the realization $\Lambda_{\zeta}$. Then,
    \begin{enumerate}
        \item For $w \in S_n \subset W$, $\rho(w)$ is a rescaled permutation matrix with non-zero entries of the form $\zeta^k$ for $k \in \Z.$ 
        For a reflection (i.e., transposition) $w= (i,j) \in S_n$ with $i < j$, $\rho(w)$ is explicitly given by 
        \[\rho(w)=
        \left({\begin{array}{ccccccccccc}
        \Id_{i-1} &    & & &  \\
           &0  &0  & \zeta^{i-j} &\\
           &0    &\Id_{j-i-1} &0   &\\
           &\zeta^{j-i}   &0   &0   & \\
           &    &   &   &\Id_{n-j}\\
           \end{array} }  \right)
        \]
         where $\Id_k$ denotes the identity matrix of size $k$ (and unspecified entries of the matrix are all 0).
        \item Conjugation with $\rho(w), \; w \in S_n$ is an honest permutation on diagonal matrices, which permutes the diagonal entries according to the permutation $w$.
        \item For a positive root $\alpha$ associated to a transposition $(i,j)$ in $S_n$ with $i<j$, let $t_{\alpha}$ be the corresponding translation. Then we have that \[ \rho(t_{\alpha})= 
          \left({\begin{array}{ccccccccccc}
        \Id_{i-1}   \\
           &\zeta^{n} \\
           &    &\Id_{j-i-1}\\
           &    &   &\zeta^{-n} & \\
           &    &   &   &\Id_{n-j}\\
           \end{array} }  \right).
        \]
        More generally, for $\lambda=\sum_{i=1}^{n-1}\lambda_i\al_{s_i} \in \Lambda_{\rt}$, we have that $\rho(t_{\lambda})$ is a diagonal matrix given by
        \[ \rho(t_{\lambda})= 
          \left({\begin{array}{ccccccccccc}
        \zeta^{n\lambda_1} \\
            &\zeta^{n(\lambda_2-\lambda_1)}\\
            &   &\zeta^{n(\lambda_3-\lambda_2)} \\
            &   &   &\ddots \\
            &   &   &   &\zeta^{n(\lambda_{n-1}-\lambda_{n-2})}\\
            &   &   &   &   &\zeta^{-n\lambda_{n-1}}
           \end{array} }  \right).
        \]
        \item Any element of $W$ acts by a rescaled permutation matrix.
    \end{enumerate}
\end{lemma}
\begin{proof}
    (1) The fact that $\rho(s_i)$ coincides with the given matrix is immediate from the definition of the $W$-action in the deformed affine realization. 
    By writing any $w\in W$ as a composition of the simple reflections $s_i$ for $1\leq i \leq n-1$, we see that $\rho(w)$ is a generalized permutation matrix as desired. 
    
    Now, observe that as an $S_n-$module (for the finite Weyl group $S_n \subset W$), we can identify $\Lambda_{\zeta}$ with the permutation representation $\Z[\zeta,\zeta^{-1}][y_1, \ldots ,y_n]$ of $S_n$ over the ring $\Z[\zeta,\zeta^{-1}]$, by setting $y_k=\zeta^{k-1}x_k$ for $1\leq k \leq n$. 
    Then we see that for a transposition $w=(i,j) \in S_n$ with $i<j,$ 
    $$\rho(w)(x_k)=\rho(w)(\zeta^{1-k}y_k)= 
    \begin{cases}
    \zeta^{1-i}y_j=\zeta^{j-i}x_j \;\text{  if $k=i$}\\
    \zeta^{1-j}y_i=\zeta^{i-j}x_i \;\text{  if $k=j$}\\
    \zeta^{1-k}y_k=x_k \hspace{18pt}\text{  otherwise}
    \end{cases},$$
   so that
    \[\rho(w)=
        \left({\begin{array}{ccccccccccc}
        \Id_{i-1} &    & & &  \\
           &0  &0  & \zeta^{i-j} &\\
           &0    &\Id_{j-i-1} &0   &\\
           &\zeta^{j-i}   &0   &0   & \\
           &    &   &   &\Id_{n-j}\\
           \end{array} }  \right)
        .\]
    
    (2) follows from (1), since we can then write $\rho(w)$ as a permutation matrix times a diagonal matrix for any $w\in S_n$.
    
    (3) We know that $s_0=t_{\al_{long}}s_{\al_{long}}$ acts via the matrix
    \[\left({\begin{array}{ccccccccccc}
             
         0  &0 &\zeta\\
         0  &\Id_{n-2} &0\\
         \zeta^{-1} &0 &0  \\
           \end{array} }  \right).
        \]
    But $s_{\al_{long}}$ acts via  the matrix 
    \[\left({\begin{array}{ccccccccccc}
             
         0  &0 &\zeta^{1-n}\\
         0  &\Id_{n-2} &0\\
         \zeta^{n-1} &0 &0  \\
           \end{array} }  \right),
        \]
    so that 
    \[\rho(t_{\al_{long}})=\rho(s_0)\rho(s_{\al_{long}})= 
    \left({\begin{array}{ccccccccccc}
             
         \zeta^{n} &0 &0\\
         0  &\Id_{n-2} &0\\
         0 &0 &\zeta^{-n}  \\
           \end{array} }  \right).
        \]
        Now we use that any positive root $\alpha$ is $S_n$-conjugate to $\alpha_{long}$ so that there is some $w \in S_n$ such that $w(\alpha_{long})=\alpha$. 
        Then $t_\alpha=t_{w(\alpha_{long})}=wt_{\alpha_{long}}w^{-1}$. 
        By part (2), we see that $\rho(t_{\al})$ is then a diagonal matrix whose diagonal entries are given by permuting the diagonal entries of $\rho(t_{\al_{long}}).$ 
        If $\alpha$ corresponds to a transposition $(i,j)$ with $i<j$, then since $\alpha_{long}$ corresponds to the transposition $(1,n)$, we may choose $w=(1,i)(n,j) \in S_n$ so that $ws_{\al_{long}}w^{-1}=(i,j)=s_{\al}$ to get the desired description of $\rho(t_{\alpha})$ using part (2) of the lemma. The general description of $\rho(t_{\lambda})$ for $\lambda \in \Lambda_{\rt}$ follows.
        
    (4) Any element of $W$ is a product of an element of $S_n$ and a translation, and a product of rescaled permutation matrices is a rescaled permutation matrix. 
\end{proof}

\begin{lemma}\label{lemma freeness zeta infinite order} Let $\Lambda_A$ be a specialization of $\Lambda_{\zeta}$ as in \cref{defn-deformed affine realization}. For any element $y \in W$, the fixed point set $\Lambda_A^y$ and the quotient $\Lambda_A/\Lambda_A^y$ are both free modules over $A$.  The rank of $\Lambda_A^y$ over $A$ agrees with the rank of $\Lambda_{\zeta}^y$ over $\Z[\zeta, \zeta^{-1}]$ so long as $\zeta \in A^{\times}$ has infinite order. \end{lemma}

\begin{proof}
By \cref{lemma-action matrix}, $y$ acts by a rescaled permutation matrix. 
From \cref{lemma rescaled permutation matrix fixed points}, we deduce the freeness of $V_A^y$ and $V_A/V_A^y$. The entries of $t$ are powers of $\zeta$, so the element $c$ for any cycle is a power of $\zeta$. By the infinite order assumption, $c = 1$ in $A$ if and only if $c=1$ in $\Z[\zeta, \zeta^{-1}]$. \end{proof}



% Now we have the following result which tells us explicitly how the $t_{\alpha} \in W$ act on $\Lambda_{\zeta}$. Let 
% $$ \rho: W \rightarrow GL_{\Z[\zeta, \zeta^{-1}]}(\Lambda_{\zeta})$$
% be the map associated to the action on $W$ on $\Lambda_{\zeta}$.
% Now that we have generalized the notion of reflection faithfulness, we can generalize Theorem \ref{thm-reff over Q(z)} in the following way:
% \begin{theorem}\label{reff thm}
%     Let $A=\Q[\zeta,\zeta^{-1}]$ and $V_A := \Lambda_{\zeta} \otimes_{\Z[\zeta,\zeta^{-1}]}A$. Then $V_A$ is a reflection faithful representation of $(W,S)$ where $W=\Tilde{S}_n, S=\{s_i\}_{i\in \Z/n\Z}$ over $A$ for any $n \geq 3$.
% \end{theorem}
% Indeed, Theorem \ref{thm-reff over Q(z)} follows from Theorem \ref{reff thm}, applying Lemma \ref{reff flat base change} to the localization map $\Q[\zeta,\zeta^{-1}]\hookrightarrow \Q(\zeta).$ We need one more lemma which makes explicit the action of reflections and affine translations on $V_A$ before proceeding to the proof of Theorem \ref{reff thm}.

\begin{theorem}\label{reff thm} 
    The realization $\Lambda_{\zeta}$ is a reflection-faithful representation of $W$. 
    Let $A$ be any $\Z[\zeta, \zeta^{-1}]$-algebra for which the unit $\zeta \in A$ has infinite order. Then $\Lambda_A$ is reflection faithful as well.
    
    % =\Q[\zeta,\zeta^{-1}]$ and $V_A := \Lambda_{\zeta} \otimes_{\Z[\zeta,\zeta^{-1}]}A$. Then $V_A$ is a reflection faithful representation of $(W,S)$ where $W=\Tilde{S}_n, S=\{s_i\}_{i\in \Z/n\Z}$ over $A$ for any $n \geq 3$.
\end{theorem}

\begin{proof}  
   %Commenting out old argument over \Q %%%%%%%%%%%
   % \comm{
   % Let $V_{\Q}$ be the specialization of $V_A$ via $$A=\Q[\zeta,\zeta^{-1}] \rightarrow \Q: \zeta \mapsto 1.$$ Note that the associated map $W \rightarrow GL_{\Q}(V_{\Q})$ factors via the quotient $S_n$ (by lemma \ref{lemma-action matrix}) and the resulting $S_n$ representation $V_{\Q}$ is known to be reflection faithful  since it coincides with the permutation representation (a.k.a. $\mathfrak{gl}_n$-realization) by the explicit formulas in \cref{lemma-action matrix}. 

   %  To show that $V_A$ is a faithful representation, suppose we have $w\in S_n$ and $\lambda \in \Lambda_{\rt}$ such that $wt_{\lambda} \in W$  acts trivially on $V_A$, and hence on $V_{\Q}$. 
   %  Since $t_{\lambda}$ acts trivially on $V_{\Q}$, we then have that $w$ acts trivially on $V_{\Q}$ so that $w=1$ (as $V_{\Q}$ is reflection faithful, in particular, faithful for $S_n$). So we now have that $t_{\lambda}$ acts trivially on $V_A$ and by the explicit description from Lemma \ref{lemma-action matrix}, we have that $\lambda =0$.}
    %%%%%%%%%%%%%%%%%%%%%%%%%%%%%%%%%%%%%%%%%%%%%%%
    Faithfulness of $V_A$ is an immediate consequence of \cref{lemma-action matrix} and \cref{lemma rescaled permutation matrix fixed points} (for $w\in S_n, \lambda \in \Lambda_{\rt}$, the only way $wt_\lambda$ acts as the identity matrix is if $w=1$ and $\lambda=0$, by the explicit description of the corresponding matrices in \cref{lemma-action matrix} and the argument in \cref{lemma rescaled permutation matrix fixed points}).

    We are left to show that $V_A$ is reflection faithful, i.e., we have to show that (b),(c) and (d) in Definition \ref{defn-reff} hold.
    \cref{lemma freeness zeta infinite order} shows us that (b) and (c) hold, so now we only have to show that (d) holds. By \cref{lemma rescaled permutation matrix fixed points} the fixed point set of any simple reflection $s$ has codimension $1$. The same is true for any reflection, as the fixed point set of $w s w^{-1}$ is $w(\Lambda_A^{s})$. 
    
    %%%%%%%%%Commenting out old argument
    % \comm{Since $A$ is a PID, for any $w \in W$, $V_A^w$ is a submodule of the finite rank free module $V_A$, so that $V_A^w$  is also free of finite rank. So (b) holds.
    
    % Now, observe that $V/V^w$ is torsion-free:
    % Indeed, if we had $v + V^w \in V/V^w$ and $a \in A$ such that $av \in V^w$, then since the action of $W$ is $A$-linear, $w(av)=aw(v)=av$, so $w(v)=v$ as $V$ is torsion free. 
    % Hence $V/V^w$ is a torsion-free finitely generated module over a PID, hence free of finite rank, implying (c).}
    %%%%%%%%%%%%%%%
    
    % Given a reflection $s \in W$, we know that $s=ws_iw^{-1}$ for some $i\in \Z/n\Z.$ Then $V^s=w(V^{s_i})$, and we have
    % $$\rk(V/V^s)=\rk(V)- \rk(V^s)= \rk(V)- \rk(V^{s_i})=\rk(V/V^{s_i})=1,$$
    % where we used the fact that $V/V^{s_i}$ is cyclic (for example, generated by $x_i$, while $\{x_i + \zeta x_{i+1}\}\cup\{x_j\}_{j\not=i, i+1} $ is a basis for $V^{s_i}$).
    
    On the other hand, suppose we have $w=w't_{\lambda} \in W$ with $w'\in S_n, \; \lambda \in \Lambda_{\rt}$ such that  $\rk_A(\Lambda_A^w)=n-1$.
    By \cref{lemma rescaled permutation matrix fixed points}, $w'$ needs to have at least $n-1$ cycles, so either $w'$ is a reflection in $S_n$, or $w'$ is the identity. If $w'$ is the identity, so that $w = t_{\lambda}$, then we need exactly $n-1$ diagonal entries in $\rho(t_{\lambda})$ (see \cref{lemma-action matrix}) to be $1$, so that all but one cycle is monodromy-free. But $\rho(t_{\lambda})$ has determinant $1$, so the last entry is also $1$, a contradiction.
    
     Thus $w'$ is a reflection. Suppose $n \ge 3$. By further conjugating $w$ (and changing the meaning of $\lambda$) we may assume that $w = s_1 t_{\lambda}$. By \cref{lemma rescaled permutation matrix fixed points}, all $n-1$ cycles of $s_1$ must be monodromy-free. In particular, all diagonal entries in $\rho(t_{\lambda})$ must be $1$ except the first two, from which one deduces that $\lambda$ is a multiple of $\al_1$.  Thus $w = s_1 t_{k \al_1}$ for some $k \in \Z$, and $w$ is a reflection. 

    We leave to the reader the analogous proof when $n=2$.
    % Hence we have shown that (d) holds as well, so that $\Lambda_A$ is indeed reflection faithful.

    %commenting out old argument %%%%%%%%%%%%%%%
    % \comm{Since $V_A/V_A^x$ is free, we have a split exact sequence 
    % $$0 \rightarrow V_A^x \rightarrow V_A \rightarrow V_A/V_A^x \rightarrow 0$$
    % which after specializing via $A \rightarrow \Q: \zeta \mapsto 1$ gives a split exact sequence
    % $$0 \rightarrow V_A^x\otimes_A \Q \rightarrow V_{\Q}=V_A\otimes_A \Q \rightarrow (V_A/V_A^x)\otimes_A \Q \rightarrow 0 .$$
    % This implies that $ V_A^x\otimes_A \Q \subset V_{\Q}^x$ so that $\dim(V_{\Q}^x)\geq \dim( V_A^x\otimes_A \Q)=n-1$ (since $V^x_A$ was free of rank n-1), and $\dim(V_{\Q}/V_\Q^x)\leq 1$.\\
    % Note that the action of $W$ on $V_\Q$ factors through $S_n$, and coincides with the permutation representation of $S_n$ over $\Q$, which is well known to be reflection faithful. Hence the action of $x=wt_{\lambda}$ on $V_\Q$ is given by the action of $w \in S_n$, so we have that $V_\Q^x=V_\Q^w$. 
    % There are two cases to consider:
    % \begin{itemize}
    %     \item $\dim(V_\Q/V_\Q^x)=0$: In this case, $w$ acts trivially on $V_\Q$, so $w=1$ (the $S_n$ action on $V_\Q$ is faithful). Hence $x=t_{\lambda}$ and by the explicit description of $t_{\lambda}$ in Lemma \ref{lemma-action matrix}, we get a contradiction to the fact that $\rk_A(V_A^{t_{\lambda}}/V_A)=\rk_A(V_A^x/V_A)=1$.
    %     \item $\dim(V_\Q/V_\Q^x)=1$: In this case, $w \in S_n$ fixes a codimension 1 subspace of $V_\Q$, and by reflection faithfulness of $V_\Q$ as a representation of $S_n$, we know that $w$ is a reflection and is $S_n$-conjugate to $s_1$, say $w=w's_1w'^{-1}$ for $w'\in S_n$. 
    %     Since the set of reflections in $W$ is stable under conjugation, it suffices to show that $y=w'^{-1}xw'=s_1t_{w'^{-1}(\lambda)}$ is a reflection.
    %     Using the explicit description from Lemma \ref{lemma-action matrix}, we see that $\rk(V_A/V_A^y)=1$ or equivalently $\rk(V_A^y)= n-1$ if and only if $w'^{-1}(\lambda)=k\alpha_1$ for some $k \in \Z$, so that $y=s_1t_{k\alpha_1}$ is indeed a reflection in $W$ (by Lemma \ref{lemma reflections}).
    % \end{itemize}}%%%%%%%%%%%%%%%%%%%%%%%%%%%
\end{proof}

\begin{lemma}\label{lemma reff for V*} Let $A$ be any $\Z[\zeta, \zeta^{-1}]$-algebra for which the unit $\zeta \in A$ has infinite order. Then $\Lambda_A^*$ is reflection faithful as well. \end{lemma}

\begin{proof} The $W$-action on $\Lambda^*$ is the same as the $W$-action on $\Lambda$ up to replacing $\zeta$ with $\zeta^{-1}$ and vice versa. 
This follows immediately by checking that if $\{x_i^*\} \subset \Lambda^*$ denotes the dual basis to $\{x_i\} \subset \Lambda$, then 
$$
s_i(x_k^*)=
\begin{cases}
 \zeta^{-1}x_{i+1}^* \text{ if k=i}\\
 \zeta x_i^* \hspace{19pt}\text{ if k= i+1} \\
 x_k^* \hspace{23 pt}\text{ otherwise}.
\end{cases}
$$
(See \cite[Proposition 2.15]{EJY1} for the formal notion of a $\zeta-$antilinear isomorphism.)\end{proof}

Having concluded the discussion of reflection faithfulness, we discuss the other critical assumption.

\begin{lemma} \label{lem:assume2} Assumption  \hyperref[assumption 2]{(2)} holds for $\Lambda_{\zeta}$, and for any specialization $\Lambda_A$. \end{lemma}

\begin{proof} Fix a finitary subset $I\subset S$, and a maximal finitary subset $S'$ containing $I$. 
Then $S'$ is of the form $S\setminus\{k\}=\hh{k}$ for some $0\le k< n.$ 
We can then identify $\Lambda_A$ with the permutation representation $A \cdot \{y_1, \ldots y_n\}$ of  $W_{S'} \cong S_n$ by choosing $y_1=x_{k+1}, y_2=\zeta^{-1} x_{k+2}, \ldots, y_n=\zeta^{-(n-1)}x_{k+n}.$
It is well-known that Assumption \hyperref[assumption 2]{(2)} holds for the permutation representation of $S_n$ even over $\Z$ (and hence after base change to $A$). This is effectively a result of Demazure \cite{Demazure} (for a modern exposition in English see \cite[Theorem 4.2]{EKLP}). For example, one can find a basis for $R$ over $R^{\hh{0}}$ by applying Demazure operators to the polynomial $y_1^{n-1} y_2^{n-2} \cdots y_{n-1}$. \end{proof}

\begin{theorem}
Assumptions \hyperref[assumption 1]{(1)} and \hyperref[assumption 2]{(2)}  from section \ref{subsec-categorification} hold for the realization $\Lambda_{A}$ so long as $A$ is an infinite field of characteristic $\ne 2$, and $\zeta$ has infinite order in $A^{\times}$. Hence Theorems \ref{thm-categorification} and \ref{thm-hom formula} hold for $\SSBim$ when defined using $\Lambda_{A}$.
        
        % $\Lambda_{\Q(\zeta)}^*=\Hom_{\Q(\zeta)}(\Lambda_{\Q(\zeta)},\Q(\zeta))$ of $(W,S)$. 
\end{theorem}
\begin{proof}
Assumption \hyperref[assumption 1]{(1)} follows from \cref{lemma reff for V*}.
Assumption \hyperref[assumption 2]{(2)} follows from \cref{lem:assume2}.
\end{proof}

\begin{remark}\label{rmk-Abe SSBim}
 Noriyuke Abe in \cite{Abe} generalizes the construction of \cite{WillSingular} by starting with weaker assumptions, including cases where $V$ is not a faithful representation of $(W,S)$ but is a faithful representation for all finitary subgroups, and allowing $\kk$ to be a complete local Noetherian integral domain.
 Let $Q$ denote the fraction field of $R$, and $Q^I$ the $W_I$-invariant subring. 
 Now we define $\SSBim_{\Abe}$  as follows: instead of just working with $(R^I,R^J)$-bimodules $M$, one has the extra data of certain $(Q^I,Q^J)$-bimodules $(M_Q^x)_{x\in W_I\backslash W/W_J}$ and isomorphisms $Q^I\otimes_{R^I}M\otimes_{R^J}Q^J\simeq \bigoplus_{x\in W_I\backslash W/W_J}M_Q^x $ which satisfy certain conditions. See \cite[section 2]{Abe} for further details of this construction. 
 It is shown then in \cite{Abe} that Theorems \ref{thm-categorification} and Theorems \ref{thm-hom formula} hold under his weaker assumptions for $\SSBim_{\Abe}$. 
 Abe is careful not to rely on the balanced assumption. We expect Abe's construction to be useful when studying specializations of $\Lambda_{\zeta}$ where $\zeta$ is a root of unity, as such specializations are no longer faithful.
\end{remark}







\subsection{Reverse Bott-Samelsons}\label{subsec reverse bott-samelsons}

The goal of this section is to prove that the 2-functor $\GS$ is essentially surjective after taking the Karoubi envelope. A brief pedagogical interlude on ordinary Bott-Samelson bimodules will set the stage. We work with a realization satisfying Assumptions \hyperref[assumption 1]{(1)} and \hyperref[assumption 2]{(2)}.

The traditional category of Bott-Samelson bimodules $\BSBim$ (these are $(R,R$)-bimodules) is the full subcategory of $\SBSBim(\emptyset, \emptyset)$ only containing objects associated to singular expressions of the form 
\begin{equation} \label{ordinaryexp} [[\emptyset \subset \{s_1\} \supset \emptyset \subset \{s_2\} \supset \cdots \subset \{s_d\} \supset \emptyset]] .\end{equation}
Thus $\BSBim$ and $\SBSBim(\emptyset, \emptyset)$ are not the same. 
However, these two monoidal categories have the same Karoubi envelope, which is not obvious, but follows from the classification of indecomposables in each category. This is stated as \cite[Theorem 24.12]{EMTW} but the proof is not explained, so let us elaborate.

\begin{theorem} \label{thm:sameKar1} The Karoubi envelopes of (the graded, additive completions of) $\BSBim$ and $\SBSBim(\emptyset, \emptyset)$ are equivalent. \end{theorem} 

\begin{proof} Since $\BSBim \subset \SBSBim(\emptyset, \emptyset)$ we need only prove that every indecomposable summand of a singular Bott-Samelson bimodule in $\SBSBim(\emptyset, \emptyset)$ is isomorphic to a summand of an ordinary Bott-Samelson bimodule. By \cref{thm-categorification}, the isomorphism classes of indecomposable objects in $\SBimod{I}{J}$ are in bijection with double cosets $p \in W_I \backslash W / W_J$, and appears as a direct summand of $\BS(IK_{\bullet})$ whenever $IK_{\bullet}$ is a (singular) reduced expression for $p$.

When $I = J = \emptyset$, $p = \{w\}$ is a singleton, and any ordinary reduced expression for $w$ gives rise to singular reduced expression as in \eqref{ordinaryexp} (though not all singular reduced expressions are of this form). Consequently, the indecomposable bimodule $B_p$ is a summand of an ordinary Bott-Samelson bimodule (it would traditionally be called $B_w$).
\end{proof}

Ultimately, the basic combinatorial fact underlying this theorem is: every double coset $p \in W_{\emptyset} \backslash W / W_{\emptyset}$ has a singular reduced expression of the form \eqref{ordinaryexp}.

Ordinary Bott-Samelson bimodules bounce between minimally singular and subminimally singular.
In this paper, we mostly deal with the other extreme. Recall the definition of a maximal double coset from \cref{defn:maximalcoset}.

\begin{defn}
Let $(W,S)$ be an affine Weyl group equipped with a realization satisfying Assumptions \hyperref[assumption 1]{(1)} and \hyperref[assumption 2]{(2)}.
The $2$-category of \emph{maximally singular Soergel bimodules}, $\mSSBim$ denote the 1-full sub-$2$-category\footnote{So the only objects of $\mSSBim$ are $\hh{s}$ for $s \in S$, and for each $s, t \in S$ we have ${}_{\hh{s}} \mSSBim_{\hh{t}} = \SBimod{\hh{s}}{\hh{t}}$.} of $\SSBim$ whose objects are the maximal finitary subsets $\hh{s}$ of $S$. 
The $2$-category of \emph{reverse Bott-Samelson bimodules} $\mSBSBim$ is the 2-full sub-$2$-category\footnote{So the 1-morphisms are restricted, but the 2-morphisms are the same.} of $\mSSBim$ containing 1-morphisms of the form
\begin{equation} \label{reverseexp} [[\hh{a_1} \supset \hh{a_1 a_2} \subset \hh{a_2} \supset \hh{a_2 a_3} \subset \cdots \supset \hh{a_{d-1} a_d} \subset \hh{a_d}]].\end{equation}
To avoid redundancy we assume that $a_j \ne a_{j+1}$ for all $1 \le j < d$.
\end{defn}


\begin{theorem}\label{thm mSBSBim vs mSSBim in type A} In affine type $A$, the Karoubi envelope of (the graded, additive completion of) $\mSBSBim$ is equivalent to $\mSSBim$. Moreover, every expression of the form \eqref{reverseexp} is a reduced expression. \end{theorem}

\begin{proof} Since $\mSBSBim \subset \mSSBim$, we need to prove that every indecomposable singular Bott-Samelson bimodule in $\SBimod{\hh{s}}{\hh{t}}$ is a summand of a reverse Bott-Samelson bimodule (for all $s, t \in S$). Just as in the proof of Theorem \ref{thm:sameKar1}, Williamson's categorification theorem reduces this problem to proving that every maximal double coset has a singular reduced expression of the form \eqref{reverseexp}. Every expression \eqref{reverseexp} is $\GS(\ula,a_d)$ for some $\ula$, whence it is reduced by \cref{thm:GSrex}. Fixing the parabolic subgroup $\hh{a_d}$ on the right, the bijections \eqref{bijectionsa} show that any maximal double coset has the form $\psi_{a_d}(\lambda)$ for some $\lambda \in \Lambda^+_{\wt}$, and \cref{thm:GSrex} provides many reduced expressions of the desired form (associated to rexes $\ula$ for $\lambda$). 
\end{proof}

\begin{remark} The Bott-Samelson bimodules for the reduced expressions from \cref{thm:GSrex} are the images under geometric Satake of tensor products of fundamental representations. The analogous statement to \cref{thm mSBSBim vs mSSBim in type A} on the other side of the Soergelified Satake equivalence is that $\Rep(\slf_n)$ is equivalent to the Karoubi envelope of $\Fund(\slf_n)$, the full subcategory monoidally generated by fundamental representations. The corresponding statement under geometric Satake states that the corresponding reverse (singular) Bott-Samelson variety gives a resolution of singularities of a Schubert variety associated to the corresponding dominant weight within the affine Grassmannian. \end{remark}

